% Section 3, part 2 of 7: transversality, isotropy, the conifold expansion, special geometry, (CYC), the cyclic system.

\begin{lem}[Transversality]
\label{lem:gau-transv}
For every basis $\mathbf x$ of $\mathcal S$ and $a,b\in\{0,\dots,3\}$, in the fraction field of $\bbK[\ell]$,
\[
\sum_{i,j}(\Conc_{\mathbf x}^{-1})_{ij}\,\Dq^ax_i\,\Dq^bx_j=(\Jet^{-1})_{ab},
\]
where $(\Jet^{-1})_{ab}=0$ for $a+b\leq2$ and $(\Jet^{-1})_{03}=u^{-1}$. In particular the transversality sums of \cref{def:gau-conc} do not depend on the basis, $\Upsilon_0=\Upsilon_1=\Upsilon_2=0$ and $\Upsilon_3=u^{-1}$.
\end{lem}

\begin{proof}
$\bbK[\ell]$ is a domain and $\det\Wr_{\mathbf x}\neq0$ (\cref{lem:gau-conc}~(iv)), so $\Conc_{\mathbf x}=\Wr_{\mathbf x}^T\Jet\Wr_{\mathbf x}$ gives $\Jet^{-1}=\Wr_{\mathbf x}\Conc_{\mathbf x}^{-1}\Wr_{\mathbf x}^T$ over the fraction field; its entry $(a,b)$ is the left side. Let $J$ be the $4\times4$ exchange matrix, $J_{ac}=\delta_{a+c,3}$. The entry $(\Jet J)_{ac}=\Jet_{a,3-c}$ vanishes for $c<a$, because then $a+3-c\geq4$ (\cref{lem:gau-conc}~(iii)); so $\Jet J$ is upper triangular with diagonal entries $\Jet_{c,3-c}=(-1)^{3-c}u$, and so is its inverse $J\Jet^{-1}$, with diagonal entries $1/\Jet_{c,3-c}$. The entry $(J\Jet^{-1})_{ca}=(\Jet^{-1})_{3-c,a}$ vanishes for $a<c$, i.e.\ $(\Jet^{-1})_{ba}=0$ for $a+b<3$, and $(\Jet^{-1})_{03}=(J\Jet^{-1})_{33}=1/\Jet_{30}=u^{-1}$. The sums $\Upsilon_k$ are the entries $(0,k)$, and $\widetilde\Upsilon_k$ is built from the bivector $\sum_{i,j}(\Conc_{\mathbf x}^{-1})_{ij}x_i\otimes x_j\in\mathcal S\otimes\mathcal S$, which does not depend on $\mathbf x$: for the basis $\mathbf xM$, $M\in GL_4(\bbQ)$, the Gram matrix is $M^T\Conc_{\mathbf x}M$.
\end{proof}

\begin{lem}[Isotropy, formal Picard--Lefschetz]
\label{lem:gau-iso}
\begin{enumerate}
\item $\sh$ commutes with $\Dq$, fixes $f_0,f_1,f_2$ and maps $y_3$ to $y_3+\varpi_1$; so $\sh(\mathcal S)=\mathcal S$.
\item $\Conc(f,\varpi_1)=0$ for every $f\in\mathcal S\cap\bbK$, and $\Conc(y_3,\varpi_1)=1$.
\item On $\mathcal S\cap\bbK$ the form $\Conc$ has rank $2$ and radical $\bbQ\varpi_1$, and a plane in $\mathcal S\cap\bbK$ is isotropic if and only if it contains $\varpi_1$. In particular $\operatorname{span}(y_0,\varpi_1)$ is isotropic for every power-series solution $y_0$; and $\Conc(f_0,f_1+\gamma'f_2)=-\frac75+2\gamma'$ vanishes only for $\gamma'=\gamma=\frac7{10}$.
\end{enumerate}
\end{lem}

\begin{proof}
(i) $\sh$ is an automorphism with inverse $\ell\mapsto\ell-1$. The maps $\Dq\circ\sh$ and $\sh\circ\Dq$ are additive and both satisfy $X(fg)=X(f)\sh(g)+\sh(f)X(g)$; they agree on $\bbK$, where $\sh$ is the identity, and on $\ell$, since $\Dq(\ell+1)=\Dq\ell=\sh(\Dq\ell)$ because $\Dq\ell\in\bbK$. By the twisted Leibniz rule they agree on every monomial $a\ell^k$, hence on $\bbK[\ell]$. Clearly $\sh f_k=f_k$, and $\sh y_3=\varpi_1(\ell+1)+\rho_3=y_3+\varpi_1$.

(ii) $\Conc(f,g)$ is a polynomial in the derivatives of $f$ and $g$ with coefficients in $\bbK$, so $\sh\Conc(f,g)=\Conc(\sh f,\sh g)$ by (i). For $f\in\mathcal S\cap\bbK$ the number $\Conc(f,y_3)\in\bbQ$ (\cref{lem:gau-conc}~(ii)) is fixed by $\sh$, so $\Conc(f,y_3)=\Conc(f,y_3+\varpi_1)=\Conc(f,y_3)+\Conc(f,\varpi_1)$, i.e.\ $\Conc(f,\varpi_1)=0$. By alternation and \cref{lem:gau-conc}~(iv), $\Conc(y_3,\varpi_1)=-\Conc(f_1,y_3)-\gamma\Conc(f_2,y_3)=1$.

(iii) By \cref{lem:gau-conc}~(iv) the Gram matrix of $\Conc$ on $(f_0,f_1,f_2)$ is $\frac75(E_{10}-E_{01})+2(E_{02}-E_{20})$, of rank $2$, with kernel spanned by $(0,1,\gamma)$, i.e.\ by $\varpi_1$; this also gives the last assertion. If a plane $\Lambda\subset\mathcal S\cap\bbK$ does not contain $\varpi_1$, then $\mathcal S\cap\bbK=\Lambda\oplus\bbQ\varpi_1$ with $\varpi_1$ in the radical, so $\Conc|_\Lambda$ has rank $2$ and $\Lambda$ is not isotropic. A plane containing $\varpi_1$ is spanned by $\varpi_1$ and some $y_0$, and it is isotropic by (ii) and alternation.
\end{proof}

\begin{lem}[The conifold expansion is a ring map]
\label{lem:gau-confexp}
Let $(\alpha,\beta)\in\bbQ^2$, $\varpi_0=f_0+\alpha\varpi_1+\beta f_2$ and $\tc=\varpi_1/\varpi_0$ (\cref{def:set-frames}), and extend $\Dq$ from $\bbK$ to $\bbQ\lser s$, $s^5=u$, by $\Dq s=-(1-s^5)/(5s^4)$. The assignment of \cref{def:set-confexp} on the generators of $\Nring$ extends to a unique ring homomorphism $\Cexp_{\alpha,\beta}\colon\Nring\to\bbQ\lser s$, and $\Cexp_{\alpha,\beta}\circ\Dq=\Dq\circ\Cexp_{\alpha,\beta}$. For the elements $C_0=c_0$, $I_1=c_0\logq+i_0$, $T=I_1/C_0$ and $C_1=\mathsf W/c_0^2$ of $\Nring$ (\cref{def:set-confexp}), which evaluate to the functions of \cref{def:set-genus0},
\[
\Cexp_{\alpha,\beta}(C_0)=\varpi_0,\qquad \Cexp_{\alpha,\beta}(I_1)=\varpi_1,\qquad \Cexp_{\alpha,\beta}(\mathsf W)=\varpi_0\Dq\varpi_1-\varpi_1\Dq\varpi_0=\varpi_0^2\,\Dq\tc ,
\]
so the conifold expansion sends $C_0$, $I_1$, $T$ and $C_1$ to $\varpi_0$, $\varpi_1$, $\tc$ and $\Dq\tc$.
\end{lem}

\begin{proof}
Write $\Cexp:=\Cexp_{\alpha,\beta}$, and recall $\Xi=-\log(1-u)\in\bbQ\fps u$ (\cref{def:set-confexp}), so that $\Dq\Xi=-1$ and $\Dq^k\Xi=0$ for $k\geq2$. By \cref{def:set-confexp}, $\Cexp(L)=s^{-1}$, $\Cexp(\logq)=-\Xi$, $\Cexp(c_j)=\Dq^j\varpi_0$ and $\Cexp(i_j)=\Dq^j\varpi_1-j\Dq^{j-1}\varpi_0+\Xi\Dq^j\varpi_0=\Dq^j(\varpi_1+\Xi\varpi_0)$, the last equality by the Leibniz rule. $\Nring$ is the localisation of a polynomial ring at $L$, $c_0$ and $\mathsf W$, so the assignment extends to a unique ring homomorphism as soon as these three go to units. $s^{-1}$ and $\varpi_0$ ($\varpi_0(0)=1$) are units, and
\[
\Cexp(\mathsf W)=\varpi_0^2+\varpi_0\bigl(\Dq\varpi_1-\varpi_0+\Xi\Dq\varpi_0\bigr)-\Dq\varpi_0\bigl(\varpi_1+\Xi\varpi_0\bigr)=\varpi_0\Dq\varpi_1-\varpi_1\Dq\varpi_0=\varpi_0^2\Dq\tc,
\]
a unit because $\Dq\tc(0)=-\tc'(0)=-1$ (\cref{def:set-frames}). Also $\Cexp(I_1)=-\Xi\varpi_0+\varpi_1+\Xi\varpi_0=\varpi_1$.

$\Dq\circ\Cexp$ and $\Cexp\circ\Dq$ are additive maps $\delta$ with $\delta(xy)=\delta(x)\Cexp(y)+\Cexp(x)\delta(y)$, so it suffices to compare them on the generators of the polynomial ring: on the inverted elements $L$, $c_0$, $\mathsf W$ such a map satisfies $\delta(x^{-1})=-\Cexp(x)^{-2}\delta(x)$. For $L$: $\Dq(s^{-1})=(s^{-6}-s^{-1})/5=\Cexp((L^6-L)/5)$; for $\logq$: $\Dq(-\Xi)=1$; for $c_0,c_1,c_2,i_0,i_1$ the images $\Dq^j\varpi_0$ and $\Dq^j(\varpi_1+\Xi\varpi_0)$ are shifted in $j$ by $\Dq$. In $\Nring$, $\Dq c_3$ is fixed by $\PFfour c_0=0$ (coefficients written through $u=L^{-5}$), and $\PFfour\varpi_0=0$ gives $\Cexp(\Dq c_3)=\Dq^4\varpi_0$. Finally $\Dq i_2$ is fixed by $\Conc(c_0,I_1)=0$, where $\Dq^3I_1=c_3\logq+3c_2+\Dq i_2$ and $\Conc$ has its coefficients in $\bbQ[u]$, $u=L^{-5}$ (\cref{def:gau-conc}); this relation is affine in $\Dq i_2$ with the unit coefficient $\Jet_{03}c_0=-L^{-5}c_0$, and by the cases treated $\Cexp(\Dq^bI_1)=\Dq^b\varpi_1$ for $b\leq2$. Its image under $\Cexp$ differs from $\Conc(\varpi_0,\varpi_1)=0$ (\cref{lem:gau-iso}~(ii)) only in the last slot, so $-u\varpi_0\bigl(\Cexp(\Dq^3I_1)-\Dq^3\varpi_1\bigr)=0$, and $\Cexp(\Dq i_2)=\Dq^3\varpi_1+\Xi\Dq^3\varpi_0-3\Dq^2\varpi_0=\Dq^3(\varpi_1+\Xi\varpi_0)=\Dq\Cexp(i_2)$. The last assertion follows from $\Cexp(I_1/c_0)=\tc$ and $\Cexp(\mathsf W/c_0^2)=\Dq\tc$.
\end{proof}

\begin{lem}[Special geometry]
\label{lem:gau-sg}
Let $y_0\in\bbK^\times$ and $t\in\bbK$ with $\Dq t\neq0$.
\begin{enumerate}
\item[(a)] For every basis $\mathbf x$ of $\mathcal S$: $\widetilde\Upsilon_1=\widetilde\Upsilon_2=0$ and $\widetilde\Upsilon_3=1/(u\,y_0^2(\Dq t)^3)$.
\item[(b)] Suppose moreover that $y_0$ and $y_1:=ty_0$ lie in $\mathcal S\cap\bbK$ and span an isotropic plane; for instance $y_1=\varpi_1$ and $y_0$ any power-series solution not proportional to $\varpi_1$ (\cref{lem:gau-iso}). Then $\mathcal S$ has a basis $(x_0,\dots,x_3)$ with $x_0=y_0$, $x_1=y_1$, $\Conc(x_0,x_3)=\Conc(x_1,x_2)=1$ and $\Conc(x_0,x_2)=\Conc(x_1,x_3)=\Conc(x_2,x_3)=0$. With $\phi_i:=x_i/y_0$ (so $\phi_0=1$, $\phi_1=t$) and $\mathrm{Yuk}:=\de_t^2\phi_2$:
\begin{enumerate}
\item $\de_t^2\phi_3=-t\,\mathrm{Yuk}$; hence $\de_t^2(y/y_0)\in\mathrm{Yuk}\cdot\operatorname{span}_\bbQ(1,t)$ for every $y\in\mathcal S$;
\item $\mathrm{Yuk}=1/(u\,y_0^2(\Dq t)^3)$.
\end{enumerate}
\end{enumerate}
\end{lem}

\begin{proof}
(a) Put $\dot t:=\Dq t$, a unit of $\bbK$, and (local notation) $\Upsilon^{\Dq}_m:=\sum_{i,j}(\Conc_{\mathbf x}^{-1})_{ij}(x_i/y_0)\Dq^m(x_j/y_0)$; $\Upsilon^{\Dq}_0=0$ because $\Conc_{\mathbf x}^{-1}$ is antisymmetric. The Leibniz rule for $x_j=y_0\cdot(x_j/y_0)$ gives $\Upsilon_k=\sum_m\binom km\,y_0\Dq^{k-m}y_0\,\Upsilon^{\Dq}_m$, so
\[
\Upsilon_1=y_0^2\Upsilon^{\Dq}_1,\quad \Upsilon_2=2y_0\Dq y_0\,\Upsilon^{\Dq}_1+y_0^2\Upsilon^{\Dq}_2,\quad \Upsilon_3=3y_0\Dq^2y_0\,\Upsilon^{\Dq}_1+3y_0\Dq y_0\,\Upsilon^{\Dq}_2+y_0^2\Upsilon^{\Dq}_3 .
\]
As operators $\Dq=\dot t\de_t$, $\Dq^2=(\Dq\dot t)\de_t+\dot t^2\de_t^2$ and $\Dq^3=(\Dq^2\dot t)\de_t+3\dot t(\Dq\dot t)\de_t^2+\dot t^3\de_t^3$, hence $\Upsilon^{\Dq}_1=\dot t\,\widetilde\Upsilon_1$, $\Upsilon^{\Dq}_2=(\Dq\dot t)\widetilde\Upsilon_1+\dot t^2\widetilde\Upsilon_2$ and $\Upsilon^{\Dq}_3=(\Dq^2\dot t)\widetilde\Upsilon_1+3\dot t(\Dq\dot t)\widetilde\Upsilon_2+\dot t^3\widetilde\Upsilon_3$. By \cref{lem:gau-transv}, $\Upsilon_1=\Upsilon_2=0$ and $\Upsilon_3=u^{-1}$; since $y_0$ and $\dot t$ are units, solving successively gives $\Upsilon^{\Dq}_1=\Upsilon^{\Dq}_2=0$, $\Upsilon^{\Dq}_3=1/(uy_0^2)$, and then (a). Isotropy is not used in (a).

(b) $y_0$ and $y_1$ are linearly independent because $t$ is not constant, so $\Lambda:=\operatorname{span}(y_0,y_1)$ is an isotropic plane. As $\Conc$ is non-degenerate on the four-dimensional space $\mathcal S$ (\cref{lem:gau-conc}~(iv)), $\Lambda^\perp$ is two-dimensional and contains $\Lambda$, so $\Lambda^\perp=\Lambda$, and the map $\mathcal S\to\Lambda^*$, $x\mapsto\Conc(\cdot,x)|_\Lambda$, is onto. Choose $x_3$ with $\Conc(y_0,x_3)=1$, $\Conc(y_1,x_3)=0$, and $x_2'$ with $\Conc(y_0,x_2')=0$, $\Conc(y_1,x_2')=1$, and put $x_2:=x_2'-\Conc(x_2',x_3)\,y_0$; then $\Conc(y_0,x_2)=0$, $\Conc(y_1,x_2)=1$ and $\Conc(x_2,x_3)=0$. The Gram matrix of $(x_0,\dots,x_3)$ is $E_{03}+E_{12}-E_{21}-E_{30}$, invertible, so this is a basis, and $\Conc_{\mathbf x}^{-1}=-\Conc_{\mathbf x}$. In this basis the definition of $\widetilde\Upsilon_k$ reads
\[
\widetilde\Upsilon_k=-\phi_0\,\de_t^k\phi_3+\phi_3\,\de_t^k\phi_0-\phi_1\,\de_t^k\phi_2+\phi_2\,\de_t^k\phi_1 .
\]
For $k\geq2$, $\de_t^k1=\de_t^kt=0$, so $\widetilde\Upsilon_k=-\de_t^k\phi_3-t\,\de_t^k\phi_2$. For $k=2$ this gives $0=-\de_t^2\phi_3-t\,\mathrm{Yuk}$, which is (i); for $k=3$ it gives $\widetilde\Upsilon_3=\de_t(t\,\mathrm{Yuk})-t\,\de_t\mathrm{Yuk}=\mathrm{Yuk}$, which is (ii) by (a). Finally, for $y=\sum_ic_ix_i$ with $c_i\in\bbQ$, $\de_t^2(y/y_0)=c_2\mathrm{Yuk}-c_3t\,\mathrm{Yuk}$.
\end{proof}

In a conifold frame, $\mathrm{Yuk}=1/(u\varpi_0^2(\Dq\tc)^3)$ (with $u=A_4$) is the conifold expansion of $\kappa^{\mathrm{LR}}_1$ of \cref{prop:set-cyclicLR}; it is normalised by the Darboux form of (b), the three-point function $\eta(H*H,H)=5\kappa^{\mathrm{LR}}_1$ carries the factor $5$ of $\eta$, and the coefficient $\kappa_1=\mathrm{Yuk}\cdot\Dq\tc$ of \cref{prop:gau-quintic} carries none.

\begin{thm}[The cyclic form of $\PFfour$, (CYC)]
\label{thm:gau-cyc}
Let $y_0\in\mathcal S\cap\bbK$ be non-zero and not proportional to $\varpi_1$, and let $\varpi_1$ be taken at any non-zero scale. Then, as operators in $\bbK[\Dq]$,
\begin{equation}
\label{eq:gau-cyc}
\PFfour=\Mcyc(y_0,\varpi_1)=\frac1{y_0}\circ\Dq\circ\frac1{\Dq t}\circ\Dq\circ\bigl(u\,y_0^2(\Dq t)^2\bigr)\circ\Dq\circ\frac1{\Dq t}\circ\Dq\circ\frac1{y_0},\qquad t=\frac{\varpi_1}{y_0}.
\end{equation}
More generally $\PFfour=\Mcyc(y_0,y_1)$ for every basis $(y_0,y_1)$ of an isotropic plane in $\mathcal S\cap\bbK$; these are the planes that contain $\varpi_1$ (\cref{lem:gau-iso}~(iii)). $\Mcyc(y_0,y_1)$ is unchanged under $y_0\mapsto\mu y_0$, $y_1\mapsto\mu'y_1$ ($\mu,\mu'\in\bbQ^\times$) and $y_1\mapsto y_1+ay_0$ ($a\in\bbQ$). In particular $\PFfour=\Mcyc(\varpi_0,\varpi_1)$ in every conifold frame of \cref{def:set-frames}.
\end{thm}

\begin{proof}
Put $t:=y_1/y_0$ and $\dot t:=\Dq t$. As operators $\Dq\circ\dot t^{-1}\circ\Dq=\dot t\,\de_t^2$, so
\[
\Mcyc(y_0,y_1)=\frac{\dot t}{y_0}\circ\de_t^2\circ\bigl(u\,y_0^2\dot t^3\bigr)\circ\de_t^2\circ\frac1{y_0}.
\]
For $y\in\mathcal S$, \cref{lem:gau-sg}~(b) gives $u\,y_0^2\dot t^3\de_t^2(y/y_0)=c_2-c_3t$ with constants $c_2,c_3$, which $\de_t^2$ kills; so $\Mcyc(y_0,y_1)$ annihilates $\mathcal S$. The coefficient of $\Dq^4$ in \eqref{eq:gau-Mcyc} is the product of the five multiplication factors,
\[
y_0^{-1}\cdot\dot t^{-1}\cdot uy_0^2\dot t^2\cdot\dot t^{-1}\cdot y_0^{-1}=u=A_4.
\] Hence $E:=\Mcyc(y_0,y_1)-\PFfour=\sum_{k\leq3}e_k\Dq^k\in\bbK[\Dq]$ annihilates $\mathcal S$, i.e.\ the row vector $(e_0,\dots,e_3)$ satisfies $(e_0,\dots,e_3)\Wr_{\mathbf f}=0$; as $\det\Wr_{\mathbf f}\neq0$ (\cref{lem:gau-conc}~(iv)), $E=0$. Every step takes place in $\bbK$ or $\bbK[\ell]$, where $y_0$ and $\dot t$ are units; the condition $y_0(0)\neq0$ is not used. Under $y_0\mapsto\mu y_0$, $y_1\mapsto\mu'y_1$ the multiplication factors of \eqref{eq:gau-Mcyc} are multiplied by $\mu^{-1},\mu\mu'^{-1},\mu'^2,\mu\mu'^{-1},\mu^{-1}$, with product $1$, and constants commute with $\Dq$; under $y_1\mapsto y_1+ay_0$, $t\mapsto t+a$ and $\dot t$ is unchanged. The plane of the first statement contains $\varpi_1$ and is isotropic by \cref{lem:gau-iso}; in a conifold frame $\varpi_0$ is a power-series solution with $\varpi_0(0)=1$, not proportional to $\varpi_1$.
\end{proof}

\begin{cor}[The cyclic system is $\PFfive$]
\label{cor:gau-cycsys}
Let $(y_0,y_1)$ be as in \cref{thm:gau-cyc}, $t=y_1/y_0$, and
\[
\kappa=(\kappa_0,\dots,\kappa_4):=\Bigl(\Dq t,\ \frac{L^5}{y_0^2(\Dq t)^2},\ \Dq t,\ y_0,\ y_0\Bigr),\qquad L^5:=u^{-1},
\]
units of $\bbK$. Let $(\cR,\Dq)$ be a differential ring extending $(\bbK,\Dq)$ and $z\in\cR$ with $\Dq z=0$.
\begin{enumerate}
\item For $Y_0,\dots,Y_4\in\cR$ the cyclic system
\begin{equation}
\label{eq:gau-system}
z\Dq Y_a=\kappa_aY_{a+1}\qquad(a\in\bbZ/5)
\end{equation}
holds if and only if $y_0Y_0$ solves $\PFfive$ and $Y_{a+1}=z\Dq Y_a/\kappa_a$ for $a=0,\dots,3$.
\item If $z$ is invertible in $\cR$, then for every $w\in\cR$ the same holds with $z\Dq$ replaced by $z\Dq+w$ in \eqref{eq:gau-system} and in $\PFfive$.
\end{enumerate}
\end{cor}

\begin{proof}
Inserting the $\kappa_a$ and comparing with \eqref{eq:gau-Mcyc} gives the identity of operators
\[
\kappa_4^{-1}\Dq\,\kappa_3^{-1}\Dq\,\kappa_2^{-1}\Dq\,\kappa_1^{-1}\Dq\,\kappa_0^{-1}\Dq\circ y_0^{-1}=y_0^{-1}\Dq\,y_0^{-1}\Dq\,\dot t^{-1}\Dq\,(uy_0^2\dot t^2)\Dq\,\dot t^{-1}\Dq\,y_0^{-1}=y_0^{-1}\Dq\circ\Mcyc(y_0,y_1),
\]
with $\dot t=\Dq t$ and $\kappa_1^{-1}=uy_0^2\dot t^2$. (i) Given the four relations $Y_{a+1}=z\Dq Y_a/\kappa_a$, $a=0,\dots,3$, the closing relation $z\Dq Y_4=\kappa_4Y_0$ reads $Y_0=z^5\kappa_4^{-1}\Dq\kappa_3^{-1}\Dq\cdots\kappa_0^{-1}\Dq\,Y_0$, because $z$ commutes with $\Dq$; by the identity and \cref{thm:gau-cyc} this is $y_0Y_0=z^5\Dq\,\PFfour(y_0Y_0)$, which is $\PFfive(y_0Y_0)=0$ by \cref{lem:set-pf5}. (That lemma is an identity of operators with coefficients in $\bbQ[q][z]$: its proof uses only $\Dq\circ q=q\circ(\Dq+1)$, which holds in $\cR$ for $q=(1-u)/5^5\in\bbK$.) The same chain of equivalences read in the other direction gives the converse. (ii) If $z$ is invertible, the map of the ring of differential operators $\cR[\Dq]$ that fixes $\cR$ and sends $\Dq$ to $\Dq+w/z$ is a ring automorphism, since $(\Dq+w/z)\circ f-f\circ(\Dq+w/z)=\Dq(f)$ for $f\in\cR$; it maps $z\Dq$ to $z\Dq+w$, and applying it to the operator identities used in (i) gives (ii).
\end{proof}

The corollary is applied in $\bbQ(\zeta)\lser s\lser z$ with $w=\zeta^iL=\zeta^is^{-1}$, to amplitudes of formal WKB solutions with exponential $\re^{u_i/z}$, $\Dq u_i=\zeta^iL$: the expanded rows (\cref{prop:gau-quintic}) and the canonical gauge applied to the model columns (\cref{thm:gau-R0}), and, through $z=s^4\zh$, to the hatted series of \cref{lem:gau-HPF}; with $w=0$ it applies in $\bbQ\lser z\lser u$ to solutions without exponential, such as the vanishing solutions of \cref{def:gau-van}. In a conifold frame $\kappa$ is the vector of \cref{prop:gau-quintic}. It does not assert that WKB solutions with rows free of logarithms exist; for the rows of \cref{def:set-rows} this is part of \cref{thm:gau-R0}.

% Section 3, part 3 of 7: the quintic system and the pulled-back model; Theorem U (statement); the commutant.

%-------------------------------------------------------------------------------
\subsection{The quintic system and the pulled-back model}
\label{sec:gau-systems}

\begin{defn}[Cyclic systems with one simple pole; regular gauges]
\label{def:gau-onepole}
Let $\Bbbk$ be a field of characteristic $0$. A \emph{cyclic system with one simple pole} over $\Bbbk\fps u$ is a system $z\de_uY=AY$ with $A=\sum_{a\in\bbZ/5}b_aE_{a,a+1}$ whose coefficients have the \emph{one-pole shape}: $b_a\in\Bbbk\fps u^\times$ for $a\neq1$ and $ub_1\in\Bbbk\fps u^\times$. Its \emph{residue} is $r:=(ub_1)(0)\in\Bbbk^\times$ and its \emph{product} is $\Bprod:=\prod_ab_a\in\Bbbk\lser u$. It is \emph{self-dual} if $b_0=b_2$ and $b_3=b_4$; for $\eta=(5\,\delta_{a+b\equiv3})_{a,b}$ the entry $(c,d)$ of $\eta A$ is $5b_{d-1}\delta_{c+d\equiv4}$ and that of $A^T\eta$ is $5b_{c-1}\delta_{c+d\equiv4}$, so self-duality is $\eta A=A^T\eta$. For a second such system $\tilde A=\sum_a\tilde b_aE_{a,a+1}$, with residue $\tilde r$, the condition of \emph{equal products} is
\[
\text{(EP)}\qquad \prod_ab_a=\prod_a\tilde b_a .
\]
The \emph{gauge equation} from $\tilde A$ to $A$ is $zP'=AP-P\tilde A$ for $P=\sum_{k\geq0}z^kP_k$, $P_k\in\mathfrak{gl}_5(\Bbbk\fps u)$; then $z\de_u(P\tilde Y)=(AP-P\tilde A)\tilde Y+P\tilde A\tilde Y=AP\tilde Y$ whenever $z\de_u\tilde Y=\tilde A\tilde Y$, so $Y=P\tilde Y$ maps solutions of the second system to solutions of the first. $\Gau(A,\tilde A)$ denotes the set of all solutions $P\in\mathfrak{gl}_5(\Bbbk\fps u)\fps z$, with no condition on $P_0$; it is a $\Bbbk\fps z$-module, since $z(P\cn)'=(AP-P\tilde A)\cn$ for $\cn\in\Bbbk\fps z$. A \emph{regular gauge} is an element of $\Gau(A,\tilde A)$ with $P_0$ invertible. $P$ is \emph{graded} if $P_k=\sum_bp_{k,b}E_{b+k,b}$ for every $k$.
\end{defn}

\begin{prop}[The quintic cyclic system]
\label{prop:gau-quintic}
Let $(\alpha,\beta)\in\bbQ^2$ give a conifold frame (\cref{def:set-frames}), $\tc=\varpi_1/\varpi_0$, and
\[
\kappa:=\Bigl(\Dq\tc,\ \frac{L^5}{\varpi_0^2(\Dq\tc)^2},\ \Dq\tc,\ \varpi_0,\ \varpi_0\Bigr),\qquad L^5=u^{-1},
\]
the vector of \cref{cor:gau-cycsys} for $(y_0,y_1)=(\varpi_0,\varpi_1)$.
\begin{enumerate}
\item $\kappa_a=\Cexp_{\alpha,\beta}(C_1\kappa^{\mathrm{LR}}_a)$ for $a\in\bbZ/5$, where $C_1\kappa^{\mathrm{LR}}=(C_1,L^5/(C_0^2C_1^2),C_1,C_0,C_0)$ (\cref{prop:set-cyclicLR}) is represented in $\Nring$ by $(\mathsf W/c_0^2,\ L^5c_0^2/\mathsf W^2,\ \mathsf W/c_0^2,\ c_0,\ c_0)$. Hence the conifold expansion carries the forward relations $(z\Dq+\zeta^iL)X_a=C_1\kappa^{\mathrm{LR}}_aX_{a+1}$ ($a=0,\dots,3$) of elements of $\Nring\lser z\otimes\bbQ(\zeta)$ to the relations $(z\Dq+\zeta^iL)\Cexp(X_a)=\kappa_a\Cexp(X_{a+1})$, $L=s^{-1}$; the closing relation is then given by \cref{cor:gau-cycsys}.
\item The cyclic system \eqref{eq:gau-system} is $z\de_uY=AY$ with $A=\sum_ab_aE_{a,a+1}$, $b_a:=-\kappa_a/(1-u)$,
\[
b=\Bigl(\tc',\ -\frac1{u\varpi_0^2(1-u)^3\tc'^2},\ \tc',\ -\frac{\varpi_0}{1-u},\ -\frac{\varpi_0}{1-u}\Bigr),
\]
a self-dual cyclic system with one simple pole over $\bbQ\fps u$, with residue $r=-1$, product $\prod_ab_a=-1/(u(1-u)^5)$, and $b_a(0)=1,1,-1,-1$ for $a=0,2,3,4$.
\item More generally, for $y_0\in\mathcal S\cap\bbK$ not proportional to $\varpi_1$ and $t=\varpi_1/y_0$, the coefficients $-\kappa_a/(1-u)$ of \cref{cor:gau-cycsys} have the one-pole shape if and only if $y_0(0)\neq0$.
\end{enumerate}
\end{prop}

\begin{proof}
(i) By \cref{lem:gau-confexp}, $\Cexp(c_0)=\varpi_0$, $\Cexp(\mathsf W/c_0^2)=\Dq\tc$ and $\Cexp(L^5)=s^{-5}=u^{-1}$; hence $\Cexp(L^5c_0^2/\mathsf W^2)=u^{-1}\varpi_0^2/(\varpi_0^4(\Dq\tc)^2)=\kappa_1$, and the other four entries are immediate. Since $\Cexp$ is a ring homomorphism commuting with $\Dq$, with $\Cexp(L)=s^{-1}$, it carries each forward relation to the corresponding relation with $\kappa_a$. By \cref{cor:gau-cycsys}~(ii), with $w=\zeta^is^{-1}$, the closing relation holds if and only if $\varpi_0\Cexp(X_0)$ solves $\PFfive$ with $z\Dq$ replaced by $z\Dq+\zeta^iL$.

(ii) As $\Dq=-(1-u)\,\dd/\dd u$, the relation $z\Dq Y_a=\kappa_aY_{a+1}$ is $z\de_uY_a=b_aY_{a+1}$. With $\Dq\tc=-(1-u)\tc'$ and $L^5=u^{-1}$ the coefficients are those displayed. In a conifold frame $\varpi_0(0)=1$ and $\tc'(0)=1$ (\cref{def:set-frames}), so $b_0,b_2,b_3,b_4$ are units with the stated values and $ub_1$ is a unit with $(ub_1)(0)=-1$. Self-duality is $b_0=b_2$, $b_3=b_4$, and $\prod_ab_a=\tc'^2\cdot\bigl(-u^{-1}\varpi_0^{-2}(1-u)^{-3}\tc'^{-2}\bigr)\cdot\varpi_0^2(1-u)^{-2}=-u^{-1}(1-u)^{-5}$.

(iii) If $y_0(0)\neq0$, then $t=u/y_0(0)+O(u^2)$ since $\varpi_1=u+O(u^2)$, so $t'(0)\neq0$, and the computation of (ii) with $(\varpi_0,\tc)$ replaced by $(y_0,t)$ shows the one-pole shape. If $y_0(0)=0$, then $-\kappa_3/(1-u)=-y_0/(1-u)$ is not a unit.
\end{proof}

\begin{prop}[Equal products and $v(u)$]
\label{prop:gau-equalprod}
Let $v\in u\bbQ\fps u$ with $v'(0)\neq0$. The model system $z\,\dd Y_b/\dd v=b^{\mathrm m}_bY_{b+1}$, $b^{\mathrm m}=(1,-v^{-1},1,-1,-1)$ (\cref{prop:set-modrows}), pulled back along $v=v(u)$ is $z\de_uY=\tilde AY$ with $\tilde A=v'(u)\sum_bb^{\mathrm m}_b(v(u))E_{b,b+1}$, i.e.
\[
\tilde b=\bigl(v',\ -v'/v,\ v',\ -v',\ -v'\bigr),\qquad \prod_a\tilde b_a=-v'^5/v .
\]
\begin{enumerate}
\item If $v=u+O(u^2)$, $\tilde A$ is a self-dual cyclic system with one simple pole over $\bbQ\fps u$, with residue $\tilde r=-1$ and $\tilde b_a(0)=1,1,-1,-1$ for $a=0,2,3,4$.
\item For $v=u+O(u^2)$, the products of $\tilde A$ and of the quintic system of \cref{prop:gau-quintic} are equal if and only if $v'^5/v=1/(u(1-u)^5)$, if and only if $v=u\,\tau(u)^{-5/4}$ (\cref{def:set-tau}).
\item (Necessity of (EP).) If $A,\tilde A$ are cyclic systems with one simple pole over $\Bbbk\fps u$ and $\Gau(A,\tilde A)\neq0$, then $\prod_ab_a=\prod_a\tilde b_a$.
\end{enumerate}
In particular, for $v=v(u)$ and every conifold frame, the quintic system and the pulled-back model are two self-dual cyclic systems with one simple pole satisfying (EP), and $b_a(0)=\tilde b_a(0)$ for $a\neq1$, $r=\tilde r=-1$.
\end{prop}

\begin{proof}
The chain rule gives $z\de_u[Y(v(u))]=v'(u)\,(z\,\dd Y/\dd v)(v(u))$, hence $\tilde A$ and $\tilde b$; and $\prod_a\tilde b_a=v'^5\cdot(-v^{-1})\cdot(-1)^2=-v'^5/v$. (i) With $v=u+O(u^2)$, $v'$ is a unit with $v'(0)=1$ and $-uv'/v$ is a unit with value $-1$ at $u=0$.

(ii) The first equivalence is $\prod_ab_a=-1/(u(1-u)^5)$ (\cref{prop:gau-quintic}~(ii)). Work in $\bbQ\lser s$, $s^5=u$, where $\dd/\dd u=(5s^4)^{-1}\dd/\dd s$ has kernel $\bbQ$. Write $v=uw$ with $w\in1+u\bbQ\fps u$; an element of $1+u\bbQ\fps u$ has a unique $n$-th root in $1+u\bbQ\fps u$ for $n\geq1$ (binomial series; and $c^n=1$ for $c=1+c_ku^k+\dots$, $c_k\neq0$, is impossible, as $c^n=1+nc_ku^k+\dots$). Since $v'(0)=1$, the equation $v'^5=v\,u^{-1}(1-u)^{-5}=w(1-u)^{-5}$ is equivalent to $v'=w^{1/5}(1-u)^{-1}$. Put $v^{1/5}:=s\,w^{1/5}$ and $v^{4/5}:=s^4w^{4/5}$; then $\frac{\dd}{\dd u}v^{4/5}=\frac45v^{-1/5}v'$, and the equation is equivalent to $\frac{\dd}{\dd u}v^{4/5}=\frac45s^{-1}(1-u)^{-1}$. On the other hand $s^4/\tau(u)=\sum_{n\geq0}\frac{4}{5n+4}u^{n+4/5}$ (\cref{def:set-tau}), whose derivative is $\frac45\sum_nu^{n-1/5}=\frac45s^{-1}(1-u)^{-1}$. So the equation is equivalent to $\frac{\dd}{\dd u}\bigl(v^{4/5}-s^4/\tau\bigr)=0$, i.e.\ to $v^{4/5}-s^4/\tau\in\bbQ$; as this difference lies in $s^4\bbQ\fps u$, it is equivalent to $w^{4/5}=\tau^{-1}$, i.e.\ $w=\tau^{-5/4}$, by the uniqueness of roots. All steps are equivalences.

(iii) Let $P\in\Gau(A,\tilde A)$ be non-zero and $n$ minimal with $P_n\neq0$. The coefficient of $z^n$ in $zP'=AP-P\tilde A$ is $0=AP_n-P_n\tilde A$, hence $A^5P_n=P_n\tilde A^5$. Now $A^5=\sum_ab_ab_{a+1}\cdots b_{a+4}E_{a,a+5}=(\prod_ab_a)\,1$, and likewise for $\tilde A$, so $(\prod_ab_a-\prod_a\tilde b_a)P_n=0$ in $\mathfrak{gl}_5(\Bbbk\lser u)$, and $P_n\neq0$ gives the claim.

The last sentence combines (i), (ii) and \cref{prop:gau-quintic}~(ii).
\end{proof}

%-------------------------------------------------------------------------------
\subsection{Theorem U}
\label{sec:gau-U}

Throughout this subsection $\Bbbk$ is a field of characteristic $0$, $A=\sum_ab_aE_{a,a+1}$ and $\tilde A=\sum_a\tilde b_aE_{a,a+1}$ are cyclic systems with one simple pole over $\Bbbk\fps u$ satisfying (EP), and a prime is $\dd/\dd u$. The following notation is local to this subsection.
\begin{itemize}
\item $\mathfrak{gl}^{(m)}:=\bigoplus_a\Bbbk\lser uE_{a,a+m}$ for $m\in\bbZ/5$; an element of $\mathfrak{gl}^{(m)}$ has \emph{shift} $m$, and $\mathfrak{gl}_5(\Bbbk\lser u)=\bigoplus_m\mathfrak{gl}^{(m)}$.
\item $\Bprod^{(m)}_a:=\tilde b_a\tilde b_{a+1}\cdots\tilde b_{a+m-1}$ for $0\leq m\leq5$ (so $\Bprod^{(0)}_a=1$ and $\Bprod^{(5)}_a=\Bprod:=\prod_a\tilde b_a$), and $\mathcal Y:=\tilde A^{-1}$.
\item $\Theta_U:=\diag(d^U_0,\dots,d^U_4)$ with $d^U_0:=1$ and $d^U_{a+1}:=d^U_a\tilde b_a/b_a$; $\operatorname{dlog}\Theta_U:=\Theta_U^{-1}\Theta_U'$; $\mathfrak d:=(\det\Theta_U/\det\Theta_U(0))^{-1/5}$.
\item $\operatorname{pr}$, $\chi^U_j$, $P^\natural$ and the degree parts $P^{[\delta]}$ are defined in \cref{lem:gau-comm,lem:gau-scal}, \cref{thm:gau-U} and \cref{lem:gau-degree}. If $\Bbbk$ contains $\zeta$, put $\mathrm{gr}:=\diag(0,1,2,3,4)$, $\zeta^{i(\mathrm{gr}+1)}:=\diag(\zeta^{i(a+1)})_a$, and define the \emph{colour rotation} $\operatorname{Rot}_iP(u,z):=\zeta^{i(\mathrm{gr}+1)}P(u,\zeta^{-i}z)\zeta^{-i(\mathrm{gr}+1)}$ on matrices and $\operatorname{Rot}_iY(u,z):=\zeta^{i(\mathrm{gr}+1)}Y(u,\zeta^{-i}z)$ on vectors.
\end{itemize}
By (EP), $d^U_4\tilde b_4/b_4=\prod_a\tilde b_a/b_a=1=d^U_0$, so $d^U_{a+1}=d^U_a\tilde b_a/b_a$ holds for every $a\in\bbZ/5$. The $d^U_a$ are units of $\Bbbk\fps u$, since $\tilde b_a/b_a$ is a unit for $a\neq1$ and $\tilde b_1/b_1=(u\tilde b_1)/(ub_1)$; hence $\det\Theta_U/\det\Theta_U(0)\in1+u\Bbbk\fps u$, $\mathfrak d\in1+u\Bbbk\fps u$ is its binomial series, and $\mathfrak d'=-\frac15\tr(\operatorname{dlog}\Theta_U)\,\mathfrak d$. Entrywise, $(A\Theta_U)_{a,a+1}=b_ad^U_{a+1}=d^U_a\tilde b_a=(\Theta_U\tilde A)_{a,a+1}$, so $A\Theta_U=\Theta_U\tilde A$.

\begin{thm}[Theorem U]
\label{thm:gau-U}
Let $A$, $\tilde A$ be cyclic systems with one simple pole over $\Bbbk\fps u$ with equal products.
\begin{enumerate}
\item[(a)] (Existence.) There is $P^\natural\in\Gau(A,\tilde A)$ with $P^\natural_0=\mathfrak d\,\Theta_U$, a diagonal matrix with unit entries.
\item[(b)] (Uniqueness.) For every $P^\natural\in\Gau(A,\tilde A)$ with $P^\natural_0=\mathfrak d\,\Theta_U$, the map $\Bbbk\fps z\to\Gau(A,\tilde A)$, $\cn\mapsto P^\natural\cn(z)$, is bijective. Consequently every solution $P$ of the gauge equation with coefficients in $\mathfrak{gl}_5(\Bbbk\fps u)$ is $P^\natural\cn(z)$ for a unique $\cn\in\Bbbk\fps z$; its $z^0$-term is $P_0=\cn(0)\,\mathfrak d\,\Theta_U$, a diagonal matrix, and $P$ is a regular gauge if and only if $\cn(0)\neq0$. If $P$ is one regular gauge, the regular gauges are exactly the $P\cn(z)$ with $\cn\in\Bbbk\fps z^\times$. For the pair $(\tilde A,\tilde A)$: $\Gau(\tilde A,\tilde A)=\Bbbk\fps z\cdot1$.
\item[(c)] (Graded gauges.) Graded regular gauges exist; if $P$ is one, the graded regular gauges are exactly the $P\cn(z)$ with $\cn\in\Bbbk\fps{z^5}^\times$.
\item[(d)] (The field.) If all $b_a$, $\tilde b_a$ lie in $\bbQ\lser u$, then $P^\natural$ can be chosen in $\mathfrak{gl}_5(\bbQ\fps u)\fps z$; for every field $\Bbbk\supseteq\bbQ$ the solutions with coefficients in $\mathfrak{gl}_5(\Bbbk\fps u)$ are then $P^\natural\Bbbk\fps z$, and $P^\natural\cn$ has rational coefficients if and only if $\cn\in\bbQ\fps z$.
\end{enumerate}
\end{thm}

Part (a) is proved after \cref{lem:gau-scal}, parts (b) and (d) after \cref{lem:gau-lead}, and part (c) is \cref{prop:gau-graded}. The proof uses only the one-pole shape, (EP) and the characteristic $0$; convergence is never used.

\begin{lem}[The commutant]
\label{lem:gau-comm}
Let $F$ be a field of characteristic $0$, $\tilde b_0,\dots,\tilde b_4\in F^\times$, $\tilde A=\sum_a\tilde b_aE_{a,a+1}$, and define $\mathfrak{gl}^{(m)}$, $\Bprod^{(m)}_a$, $\Bprod$ and $\mathcal Y$ over $F$ as above.
\begin{enumerate}
\item[(a)] $\tilde A^m=\sum_a\Bprod^{(m)}_aE_{a,a+m}$ for $0\leq m\leq5$, so $\tilde A^5=\Bprod\,1$; $\mathcal Y=\sum_a\tilde b_a^{-1}E_{a+1,a}$, $\mathcal Y^j=\sum_a(\Bprod^{(j)}_a)^{-1}E_{a+j,a}$ for $0\leq j\leq5$, and $\tilde A^m=\Bprod\,\mathcal Y^{5-m}$ for $1\leq m\leq5$; $\det(x1-\tilde A)=x^5-\Bprod$.
\item[(b)] For $X=\sum_a\xi_aE_{a,a+m}\in\mathfrak{gl}^{(m)}$, $[\tilde A,X]=\sum_a(\tilde b_a\xi_{a+1}-\xi_a\tilde b_{a+m})E_{a,a+m+1}\in\mathfrak{gl}^{(m+1)}$.
\item[(c)] For every $m$, $\ker(\operatorname{ad}\tilde A)\cap\mathfrak{gl}^{(m)}=F\tilde A^m$; hence $\ker(\operatorname{ad}\tilde A)=F[\tilde A]=F[\mathcal Y]=\bigoplus_{j=0}^4F\mathcal Y^j$, of dimension $5$.
\item[(d)] For $m\in\{1,\dots,5\}$ and $X=\sum_a\xi_aE_{a,a+m}\in\mathfrak{gl}^{(m)}$ put $\lambda_m(X):=\sum_a\xi_a/\Bprod^{(m)}_a$ (a local letter). Then $\lambda_m(\tilde A^m)=5$, $\operatorname{ad}\tilde A(\mathfrak{gl}^{(m-1)})=\ker\lambda_m$, $\mathfrak{gl}^{(m)}=F\tilde A^m\oplus\ker\lambda_m$, and $\mathfrak{gl}_5(F)=F[\mathcal Y]\oplus\operatorname{im}(\operatorname{ad}\tilde A)$.
\item[(e)] The projection $\operatorname{pr}$ onto $F[\mathcal Y]$ along $\operatorname{im}(\operatorname{ad}\tilde A)$ is $F$-linear and preserves the shift; for $X\in\mathfrak{gl}^{(m)}$, $1\leq m\leq5$,
\[
\operatorname{pr}(X)=\tfrac15\lambda_m(X)\,\tilde A^m=\tfrac15\Bigl[\sum_a\xi_a\Bprod^{(5-m)}_{a+m}\Bigr]\mathcal Y^{5-m}.
\]
Equivalently, $\operatorname{pr}(X)$ is the unique element of $F[\mathcal Y]$ with $\tr(\operatorname{pr}(X)Z)=\tr(XZ)$ for all $Z\in F[\mathcal Y]$.
\item[(f)] Let $F=\Bbbk\lser u$ and let $\tilde A$ have the one-pole shape. Then $\mathcal Y\in\mathfrak{gl}_5(\Bbbk\fps u)$; an element $\sum_{j=0}^4w_j\mathcal Y^j$ ($w_j\in\Bbbk\lser u$) lies in $\mathfrak{gl}_5(\Bbbk\fps u)$ if and only if every $w_j$ lies in $\Bbbk\fps u$; the commutant of $\tilde A$ in $\mathfrak{gl}_5(\Bbbk\fps u)$ is the ring $\Bbbk\fps u[\mathcal Y]:=\bigoplus_{j=0}^4\Bbbk\fps u\,\mathcal Y^j$; and for $X\in\mathfrak{gl}_5(\Bbbk\fps u)$ the coefficient of $\mathcal Y^j$ in $\operatorname{pr}(X)$ lies in $u^{-1}\Bbbk\fps u$ for $j=1,\dots,4$ and in $\Bbbk\fps u$ for $j=0$.
\item[(g)] (Not used.) $x^5-\Bprod$ is separable, so $\tilde A$ is regular semisimple; in the situation of (f), $x^5-\Bprod$ is irreducible over $\Bbbk\lser u$, so $F[\mathcal Y]$ is a field.
\end{enumerate}
\end{lem}

\begin{proof}
(a) By induction, $\tilde A^{m+1}=\sum_a\Bprod^{(m)}_a\tilde b_{a+m}E_{a,a+m+1}$ and $\Bprod^{(m)}_a\tilde b_{a+m}=\Bprod^{(m+1)}_a$; for $m=5$, $E_{a,a+5}=E_{aa}$. $\tilde A\sum_c\tilde b_c^{-1}E_{c+1,c}=\sum_aE_{aa}$ gives $\mathcal Y$, and $\mathcal Y^j$ follows in the same way. Reindexing, $\mathcal Y^{5-m}=\sum_a(\Bprod^{(5-m)}_{a+m})^{-1}E_{a,a+m}$, and $\Bprod^{(m)}_a\Bprod^{(5-m)}_{a+m}=\Bprod$, so $\mathcal Y^{5-m}=\Bprod^{-1}\tilde A^m$. In the Leibniz expansion of $\det(x1-\tilde A)$ a permutation $\pi$ contributes only if $\pi(a)\in\{a,a+1\}$ for all $a$; if $\pi(a)=a+1$ for one $a$, then $\pi(a+1)\neq a+1$, so $\pi(a+1)=a+2$, and so on: $\pi$ is the $5$-cycle $a\mapsto a+1$, which is even and contributes $\prod_a(-\tilde b_a)=-\Bprod$; otherwise $\pi$ is the identity.

(b) $\tilde AX=\sum_a\tilde b_a\xi_{a+1}E_{a,a+1+m}$ and $X\tilde A=\sum_a\xi_a\tilde b_{a+m}E_{a,a+m+1}$.

(c) By (b), $\operatorname{ad}\tilde A$ maps $\mathfrak{gl}^{(m)}$ to $\mathfrak{gl}^{(m+1)}$ and the sum of the $\mathfrak{gl}^{(m)}$ is direct, so $X$ commutes with $\tilde A$ if and only if each shift component does. For $X\in\mathfrak{gl}^{(m)}$ this means $\xi_{a+1}=\xi_a\tilde b_{a+m}/\tilde b_a$ for all $a$: a solution is determined by $\xi_0$, and $\xi_a=\Bprod^{(m)}_a$ is a non-zero solution, since $\Bprod^{(m)}_{a+1}\tilde b_a=\Bprod^{(m)}_a\tilde b_{a+m}$. So $\ker(\operatorname{ad}\tilde A)\cap\mathfrak{gl}^{(m)}=F\tilde A^m$ and $\ker(\operatorname{ad}\tilde A)=\bigoplus_{m=1}^5F\tilde A^m=F[\tilde A]$. By (a), $\tilde A^m=\Bprod\mathcal Y^{5-m}$ and $\mathcal Y=\Bprod^{-1}\tilde A^4$, so $F[\tilde A]=F[\mathcal Y]=\bigoplus_{j=0}^4F\mathcal Y^j$.

(d) For $Z\in\ker(\operatorname{ad}\tilde A)$ and any $Y$, $\tr(Z[\tilde A,Y])=\tr(\tilde AZY)-\tr(Z Y\tilde A)=0$. Take $Z=\tilde A^{5-m}$ and $Y\in\mathfrak{gl}^{(m-1)}$; for $X\in\mathfrak{gl}^{(m)}$, (a) gives $X\tilde A^{5-m}=\sum_a\xi_a\Bprod^{(5-m)}_{a+m}E_{aa}$, so $\tr(X\tilde A^{5-m})=\sum_a\xi_a\Bprod/\Bprod^{(m)}_a=\Bprod\,\lambda_m(X)$. Hence $\lambda_m$ vanishes on $\operatorname{ad}\tilde A(\mathfrak{gl}^{(m-1)})$. This image has dimension $5-1=4$ by (c), and $\lambda_m\neq0$ because $\lambda_m(\tilde A^m)=\sum_a\Bprod^{(m)}_a/\Bprod^{(m)}_a=5\neq0$; so the image is $\ker\lambda_m$, and $\mathfrak{gl}^{(m)}=F\tilde A^m\oplus\ker\lambda_m$. By (b), $\operatorname{im}(\operatorname{ad}\tilde A)=\bigoplus_m\operatorname{ad}\tilde A(\mathfrak{gl}^{(m-1)})$, and summing over $m$ gives $\mathfrak{gl}_5(F)=F[\tilde A]\oplus\operatorname{im}(\operatorname{ad}\tilde A)$.

(e) $X=\frac15\lambda_m(X)\tilde A^m+\bigl(X-\frac15\lambda_m(X)\tilde A^m\bigr)$ with the second term in $\ker\lambda_m$; the second form uses $\tilde A^m=\Bprod\mathcal Y^{5-m}$ and $\Bprod/\Bprod^{(m)}_a=\Bprod^{(5-m)}_{a+m}$. $X-\operatorname{pr}(X)$ lies in $\operatorname{im}(\operatorname{ad}\tilde A)$, which is orthogonal to $F[\mathcal Y]$ under $(X,Z)\mapsto\tr(XZ)$ by the first line of (d). The trace form is non-degenerate on $F[\tilde A]$: $\tilde A^n$ has shift $n$, so $\tr\tilde A^n=0$ unless $5\mid n$, $\tr1=5$ and $\tr\tilde A^5=5\Bprod$; in the basis $1,\tilde A,\dots,\tilde A^4$ the Gram matrix has $5$ at $(0,0)$, $5\Bprod$ at $(i,5-i)$ for $i=1,\dots,4$ and $0$ elsewhere, with determinant $5(5\Bprod)^4\neq0$. This gives the uniqueness.

(f) The entries of $\mathcal Y$ are the $\tilde b_a^{-1}$, which are integral ($\tilde b_1^{-1}=u(u\tilde b_1)^{-1}$). By (a), $\mathcal Y^j$ is supported on the positions $(a+j,a)$; for $j=0,\dots,4$ these five sets partition the $25$ positions, so the entry $(a+j,a)$ of $\sum_jw_j\mathcal Y^j$ is $w_j(\Bprod^{(j)}_a)^{-1}$. For $j\leq4$ the window $\{a,\dots,a+j-1\}$ has at most four elements and contains the pole index $1$ at most once, so $\Bprod^{(j)}_a$ is a unit or $u^{-1}$ times a unit and $(\Bprod^{(j)}_a)^{-1}$ is integral; for $a=2$ the window lies in $\{2,3,4,0\}$ and $(\Bprod^{(j)}_2)^{-1}$ is a unit. Hence the element is integral if and only if every $w_j$ is. $\Bbbk\fps u[\mathcal Y]$ is a ring because $\mathcal Y^5=\Bprod^{-1}1$ and $\Bprod^{-1}\in u\Bbbk\fps u$; by (c) it is the integral commutant. For $X\in\mathfrak{gl}^{(m)}\cap\mathfrak{gl}_5(\Bbbk\fps u)$, $1\leq m\leq4$, the coefficient of $\mathcal Y^{5-m}$ in $\operatorname{pr}(X)$ is $\frac15\sum_a\xi_a\Bprod^{(5-m)}_{a+m}$, a sum of products of an integral $\xi_a$ with at most four consecutive weights, at most one of them $\tilde b_1$, so it lies in $u^{-1}\Bbbk\fps u$; for $m=5$ it is $\frac15\tr X$, which is integral.

(g) $\gcd(x^5-\Bprod,5x^4)=1$ since $\Bprod\neq0$, so $\tilde A$ has five distinct eigenvalues over an algebraic closure. In the situation of (f) write $\Bprod=\Bprod^\circ/u$ with $\Bprod^\circ:=u\Bprod\in\Bbbk\fps u^\times$ (a local letter); with $x=1/X$, $x^5-\Bprod=-\Bprod^\circ u^{-1}X^{-5}\bigl(X^5-u(\Bprod^\circ)^{-1}\bigr)$, and $X^5-u(\Bprod^\circ)^{-1}$ is an Eisenstein polynomial at the prime $u$ of $\Bbbk\fps u$.
\end{proof}

% Section 3, part 4 of 7: lifting, the scalar equation, proof of Theorem U, the leading term, the Z/5-degree, graded gauges.

\begin{lem}[Lifting]
\label{lem:gau-lift}
Let $\tilde A$ have the one-pole shape, let $X\in\mathfrak{gl}_5(\Bbbk\fps u)$ have shift $m+1$ with $0\leq m\leq4$, and suppose $\operatorname{pr}(X)=0$. Then there is $\mathcal Z\in\mathfrak{gl}_5(\Bbbk\fps u)$ of shift $m$ with $[\tilde A,\mathcal Z]=X$, and every solution $\mathcal Z'\in\mathfrak{gl}_5(\Bbbk\fps u)$ of $[\tilde A,\mathcal Z']=X$ is $\mathcal Z+w$ with $w\in\Bbbk\fps u[\mathcal Y]$. If $\tilde A$ and $X$ have rational coefficients, $\mathcal Z$ can be chosen rational.
\end{lem}

\begin{proof}
Write $X=\sum_a\xi_aE_{a,a+m+1}$ and look for $\mathcal Z=\sum_a\rho_a\Bprod^{(m)}_aE_{a,a+m}$ with $\rho_a\in\Bbbk\lser u$ (local letters). By \cref{lem:gau-comm}~(b) and $\tilde b_a\Bprod^{(m)}_{a+1}=\Bprod^{(m)}_a\tilde b_{a+m}=\Bprod^{(m+1)}_a$,
\[
[\tilde A,\mathcal Z]=\sum_a(\rho_{a+1}-\rho_a)\Bprod^{(m+1)}_aE_{a,a+m+1},
\]
so the equation is $\rho_{a+1}-\rho_a=\sigma_a:=\xi_a/\Bprod^{(m+1)}_a$ for $a\in\bbZ/5$. It is solvable if and only if $\sum_a\sigma_a=\lambda_{m+1}(X)=0$, i.e.\ if and only if $\operatorname{pr}(X)=0$ (\cref{lem:gau-comm}~(e)). Each $\sigma_a$ is integral, and $\sigma_a(0)=0$ whenever the window $\{a,\dots,a+m\}$ contains the pole index $1$, because then $\Bprod^{(m+1)}_a$ is $u^{-1}$ times a unit; for $m+1=5$ every window contains $1$. Put $\rho_1:=0$ and $\rho_{a+1}:=\rho_a+\sigma_a$ for $a=1,2,3,4$; the relation for $a=0$ then holds because $\sum_a\sigma_a=0$. All $\rho_a$ are integral. The entry $\rho_a\Bprod^{(m)}_a$ is integral if the window $\{a,\dots,a+m-1\}$ avoids $1$, and if it contains $1$ (that is, for the $m$ indices $a=2-m,\dots,1$) it is integral provided $\rho_a(0)=0$. For $a=1$ this holds by the choice $\rho_1=0$; for $a=2-m,\dots,0$ one has $\rho_a=\rho_1-\sum_{c=a}^0\sigma_c$, and every window $\{c,\dots,c+m\}$ with $2-m\leq c\leq0$ contains $1$, so $\rho_a(0)=0$. Hence $\mathcal Z$ is integral. Two integral solutions differ by an element of $\ker(\operatorname{ad}\tilde A)\cap\mathfrak{gl}_5(\Bbbk\fps u)=\Bbbk\fps u[\mathcal Y]$ (\cref{lem:gau-comm}~(c), (f)). The construction uses only field operations on the data.
\end{proof}

\begin{lem}[The scalar equation]
\label{lem:gau-scal}
Let $A$, $\tilde A$ have the one-pole shape and satisfy (EP).
\begin{enumerate}
\item[(a)] For $w=\sum_{j=0}^4w_j\mathcal Y^j$ with $w_j\in\Bbbk\lser u$,
\[
\operatorname{pr}\bigl(w'+\operatorname{dlog}\Theta_U\,w\bigr)=\sum_{j=0}^4\bigl(w_j'+\chi^U_jw_j\bigr)\mathcal Y^j,\qquad \chi^U_j:=-\frac j5\,\frac{\Bprod'}{\Bprod}+\frac15\tr(\operatorname{dlog}\Theta_U).
\]
In particular $\operatorname{pr}((\mathcal Y^j)')=-\frac j5(\Bprod'/\Bprod)\mathcal Y^j$ and $\operatorname{pr}(\operatorname{dlog}\Theta_U\,\mathcal Y^j)=\frac15\tr(\operatorname{dlog}\Theta_U)\,\mathcal Y^j$. Moreover $\chi^U_j=\frac j{5u}+\chi^\circ_j$ with $\chi^\circ_j:=-\frac j5(\Bprod^\circ)'/\Bprod^\circ+\frac15\tr(\operatorname{dlog}\Theta_U)\in\Bbbk\fps u$, where $\Bprod^\circ:=u\Bprod\in\Bbbk\fps u^\times$.
\item[(b)] For $Z\in\mathfrak{gl}_5(\Bbbk\fps u)$ the equation $\operatorname{pr}(w'+\operatorname{dlog}\Theta_U\,w)=-\operatorname{pr}(Z)$ has solutions $w\in\Bbbk\fps u[\mathcal Y]$; their components $w_1,\dots,w_4$ are unique, and $w_0$ is unique up to $w_0\mapsto w_0+\cn\,\mathfrak d$, $\cn\in\Bbbk$.
\item[(c)] For $Z=0$ the solutions in $\Bbbk\fps u[\mathcal Y]$ are exactly the $\cn\,\mathfrak d\,1$, $\cn\in\Bbbk$.
\item[(d)] If the data are rational, so are the solutions for $\cn\in\bbQ$; $\mathfrak d\in1+u\bbQ\fps u$.
\end{enumerate}
\end{lem}

\begin{proof}
(a) $(\mathcal Y^j)'=\sum_a((\Bprod^{(j)}_a)^{-1})'E_{a+j,a}$ has shift $5-j$; written as $\sum_{a'}\xi_{a'}E_{a',a'+5-j}$ it has $a'=a+j$ and $\xi_{a'}=((\Bprod^{(j)}_{a'-j})^{-1})'$. By \cref{lem:gau-comm}~(e) with $m=5-j$ (and $m=5$ for $j=0$),
\[
\operatorname{pr}((\mathcal Y^j)')=\frac15\Bigl[\sum_a\bigl((\Bprod^{(j)}_a)^{-1}\bigr)'\Bprod^{(j)}_a\Bigr]\mathcal Y^j=-\frac15\Bigl[\sum_a\frac{(\Bprod^{(j)}_a)'}{\Bprod^{(j)}_a}\Bigr]\mathcal Y^j .
\]
Here $(\Bprod^{(j)}_a)'/\Bprod^{(j)}_a=\sum_{i=0}^{j-1}\tilde b_{a+i}'/\tilde b_{a+i}$, and as $a$ runs over $\bbZ/5$ each index $c$ occurs for exactly the $j$ pairs $(a,i)=(c-i,i)$; so the sum is $j\sum_c\tilde b_c'/\tilde b_c=j\Bprod'/\Bprod$. Next, $\operatorname{dlog}\Theta_U\,\mathcal Y^j=\sum_a\frac{(d^U_{a+j})'}{d^U_{a+j}}(\Bprod^{(j)}_a)^{-1}E_{a+j,a}$, so now $\xi_{a'}=\frac{(d^U_{a'})'}{d^U_{a'}}(\Bprod^{(j)}_{a'-j})^{-1}$, and \cref{lem:gau-comm}~(e) gives $\frac15\bigl[\sum_{a'}(d^U_{a'})'/d^U_{a'}\bigr]\mathcal Y^j=\frac15\tr(\operatorname{dlog}\Theta_U)\mathcal Y^j$. As $\operatorname{pr}$ is $\Bbbk\lser u$-linear and fixes $\Bbbk\lser u[\mathcal Y]$, $\operatorname{pr}(w'+\operatorname{dlog}\Theta_Uw)=\sum_j\operatorname{pr}\bigl(w_j'\mathcal Y^j+w_j(\mathcal Y^j)'+w_j\operatorname{dlog}\Theta_U\mathcal Y^j\bigr)=\sum_j(w_j'+\chi^U_jw_j)\mathcal Y^j$. Finally $\Bprod=\Bprod^\circ/u$ with $\Bprod^\circ$ a unit by the one-pole shape, so $\Bprod'/\Bprod=(\Bprod^\circ)'/\Bprod^\circ-u^{-1}$, and $\tr(\operatorname{dlog}\Theta_U)=\sum_c(d^U_c)'/d^U_c$ is integral because the $d^U_c$ are units.

(b) By \cref{lem:gau-comm}~(f), $\operatorname{pr}(Z)=\sum_j\phi^Z_j\mathcal Y^j$ with $u\phi^Z_j\in\Bbbk\fps u$ for $j\geq1$ and $\phi^Z_0\in\Bbbk\fps u$ (local letters). By (a) and the linear independence of the $\mathcal Y^j$ the equation splits into $w_j'+\chi^U_jw_j=-\phi^Z_j$, $j=0,\dots,4$. For $j=1,\dots,4$ multiply by $u$: $uw_j'+(\frac j5+u\chi^\circ_j)w_j=-u\phi^Z_j$. For $w_j=\sum_nw_{j,n}u^n$ the coefficient of $u^n$ reads
\[
\Bigl(n+\frac j5\Bigr)w_{j,n}=-\coef{u^n}(u\phi^Z_j)-\sum_{l=0}^{n-1}\coef{u^{n-1-l}}\chi^\circ_j\cdot w_{j,l},
\]
and $n+j/5$ is a non-zero rational number for every $n\geq0$: so there is exactly one solution $w_j\in\Bbbk\fps u$. For $j=0$, $\chi^U_0=\frac15\tr(\operatorname{dlog}\Theta_U)=-\mathfrak d'/\mathfrak d$, so the equation is $\mathfrak d\,(w_0/\mathfrak d)'=-\phi^Z_0$; its solutions in $\Bbbk\fps u$ are $w_0=\mathfrak d\bigl(\cn-\int_0^u\phi^Z_0/\mathfrak d\,\dd u\bigr)$, $\cn\in\Bbbk$, where $\int_0^u$ is the formal antiderivative without constant term, and the difference $e$ of two solutions satisfies $(e/\mathfrak d)'=0$, so $e\in\Bbbk\mathfrak d$.

(c) is (b) with $\phi^Z=0$: the recursion gives $w_j=0$ for $j\geq1$, and $w_0=\cn\mathfrak d$. (d) The pivots $n+j/5$ are rational, $\det\Theta_U/\det\Theta_U(0)$ lies in $1+u\bbQ\fps u$ and its binomial series $\mathfrak d$ has rational coefficients.
\end{proof}

\begin{proof}[Proof of \cref{thm:gau-U}~(a)]
\emph{Reduction.} Write $P=\Theta_U\check P$; $\check P$ (a local letter) has coefficients in $\mathfrak{gl}_5(\Bbbk\fps u)$ if and only if $P$ has, since $\Theta_U$ is invertible over $\Bbbk\fps u$. As $A\Theta_U=\Theta_U\tilde A$, we have $AP-P\tilde A=\Theta_U(\tilde A\check P-\check P\tilde A)$ and $zP'=z\Theta_U'\check P+z\Theta_U\check P'$, so the gauge equation is equivalent to $z\check P'=[\tilde A,\check P]-z\operatorname{dlog}\Theta_U\,\check P$, i.e.\ to
\begin{equation}
\label{eq:gau-reduced}
[\tilde A,\check P_0]=0,\qquad [\tilde A,\check P_{k+1}]=\check P_k'+\operatorname{dlog}\Theta_U\,\check P_k\quad(k\geq0).
\end{equation}
\emph{Construction.} We construct $\check P_k=\mathcal Z_k+w^{(k)}$ with $\mathcal Z_k\in\mathfrak{gl}_5(\Bbbk\fps u)$ and $w^{(k)}\in\Bbbk\fps u[\mathcal Y]$ such that, for every $k\geq0$, (I$_k$) $[\tilde A,\check P_k]=\check P_{k-1}'+\operatorname{dlog}\Theta_U\check P_{k-1}$ (with $\check P_{-1}:=0$) and (S$_k$) $\operatorname{pr}(\check P_k'+\operatorname{dlog}\Theta_U\check P_k)=0$. For $k=0$ put $\mathcal Z_0:=0$ and $w^{(0)}:=\mathfrak d\,1$: (I$_0$) holds because $w^{(0)}$ commutes with $\tilde A$, and (S$_0$) is \cref{lem:gau-scal}~(c). Given $\check P_k$, the matrix $X_{k+1}:=\check P_k'+\operatorname{dlog}\Theta_U\check P_k$ is integral with $\operatorname{pr}(X_{k+1})=0$ by (S$_k$); since $\operatorname{pr}$ preserves the shift, \cref{lem:gau-lift} applied to each shift component gives an integral $\mathcal Z_{k+1}$ with $[\tilde A,\mathcal Z_{k+1}]=X_{k+1}$. Choose $w^{(k+1)}\in\Bbbk\fps u[\mathcal Y]$ by \cref{lem:gau-scal}~(b) with $Z:=\mathcal Z_{k+1}'+\operatorname{dlog}\Theta_U\mathcal Z_{k+1}$ and any value of the free constant. Then (I$_{k+1}$) holds since $[\tilde A,w^{(k+1)}]=0$, and (S$_{k+1}$) holds since $\operatorname{pr}(\check P_{k+1}'+\operatorname{dlog}\Theta_U\check P_{k+1})=\operatorname{pr}(Z)+\operatorname{pr}\bigl((w^{(k+1)})'+\operatorname{dlog}\Theta_Uw^{(k+1)}\bigr)=0$. By \eqref{eq:gau-reduced}, $P^\natural:=\Theta_U\sum_kz^k\check P_k$ lies in $\Gau(A,\tilde A)$, and $P^\natural_0=\mathfrak d\,\Theta_U$. The only choices are the free constant of \cref{lem:gau-scal}~(b) at each order and the normalisation $\rho_1=0$ inside \cref{lem:gau-lift}, which the next choice of $w^{(k+1)}$ absorbs.
\end{proof}

\begin{lem}[The leading term of a solution of the gauge equation]
\label{lem:gau-lead}
Let $A$, $\tilde A$ have the one-pole shape and satisfy (EP), let $P\in\Gau(A,\tilde A)$ and $n\geq0$ with $P_0=\dots=P_{n-1}=0$. Then $P_n=\cn\,\mathfrak d\,\Theta_U$ for a unique $\cn\in\Bbbk$, namely $\cn=(P_n)_{00}(0)$.
\end{lem}

\begin{proof}
Put $\check P:=\Theta_U^{-1}P$, so $\check P_l=0$ for $l<n$. The coefficient of $z^n$ in the reduced equation \eqref{eq:gau-reduced} is $[\tilde A,\check P_n]=0$ (for $n=0$ this is the first equation), so $\check P_n\in\ker(\operatorname{ad}\tilde A)\cap\mathfrak{gl}_5(\Bbbk\fps u)=\Bbbk\fps u[\mathcal Y]$ (\cref{lem:gau-comm}~(c), (f)). The coefficient of $z^{n+1}$ says that $\check P_n'+\operatorname{dlog}\Theta_U\check P_n=[\tilde A,\check P_{n+1}]$ lies in $\operatorname{im}(\operatorname{ad}\tilde A)$, so $\operatorname{pr}(\check P_n'+\operatorname{dlog}\Theta_U\check P_n)=0$. By \cref{lem:gau-scal}~(c), $\check P_n=\cn\,\mathfrak d\,1$. As $d^U_0=1$ and $\mathfrak d(0)=1$, $\cn=(P_n)_{00}(0)$.
\end{proof}

\begin{proof}[Proof of \cref{thm:gau-U}~(b) and (d)]
(b) Let $P^\natural\in\Gau(A,\tilde A)$ with $P^\natural_0=\mathfrak d\,\Theta_U$; the gauge of (a) is one. \emph{Injectivity:} if $\cn\neq0$ has order $n$ in $z$, then $\coef{z^n}(P^\natural\cn)=\cn_nP^\natural_0=\cn_n\mathfrak d\Theta_U\neq0$. \emph{Surjectivity:} let $P\in\Gau(A,\tilde A)$. We define $\cn_0,\cn_1,\dots\in\Bbbk$ recursively. Suppose $R^{(n)}:=P-P^\natural(\cn_0+\cn_1z+\dots+\cn_{n-1}z^{n-1})$ is $0$ modulo $z^n$ (for $n=0$ this is empty); $R^{(n)}\in\Gau(A,\tilde A)$ because $\Gau(A,\tilde A)$ is a $\Bbbk\fps z$-module. By \cref{lem:gau-lead}, $\coef{z^n}R^{(n)}=\cn_n\mathfrak d\Theta_U$ for a unique $\cn_n\in\Bbbk$. Then $R^{(n+1)}=R^{(n)}-\cn_nz^nP^\natural$ lies in $\Gau(A,\tilde A)$ and $\coef{z^n}R^{(n+1)}=\cn_n\mathfrak d\Theta_U-\cn_nP^\natural_0=0$, so $R^{(n+1)}\equiv0$ modulo $z^{n+1}$. Put $\cn:=\sum_n\cn_nz^n$. For every $l$ and every $n>l$, $\coef{z^l}(P-P^\natural\cn)=\coef{z^l}R^{(n)}-\coef{z^l}P^\natural(\cn-\sum_{m<n}\cn_mz^m)=0$. Hence $P=P^\natural\cn$. Its $z^0$-term is $\cn(0)\mathfrak d\Theta_U$, invertible if and only if $\cn(0)\neq0$. If $P=P^\natural\cn_P$ is a regular gauge, every regular gauge $P^\natural\cn'$ equals $P\cdot(\cn_P^{-1}\cn')$ with $\cn_P^{-1}\cn'\in\Bbbk\fps z^\times$, and conversely. For the pair $(\tilde A,\tilde A)$ one has $\Theta_U=1$ and $\mathfrak d=1$, and $1\in\Gau(\tilde A,\tilde A)$; so $\Gau(\tilde A,\tilde A)=\Bbbk\fps z\cdot1$.

(d) Every object used in the proof of (a) is obtained from the $b_a$, $\tilde b_a$ by field operations in $\bbQ\lser u$, the derivative and the integration $\int_0^u$ of power series: $\Theta_U$, $\operatorname{dlog}\Theta_U$, $\mathcal Y$, $\Bprod^{(m)}_a$, $\Bprod^\circ$, $\operatorname{pr}$ (\cref{lem:gau-comm}~(e)), $\mathfrak d$, the lift of \cref{lem:gau-lift} and the solutions of \cref{lem:gau-scal}~(b) (\cref{lem:gau-scal}~(d)). With every free constant chosen in $\bbQ$, the construction therefore runs inside $\mathfrak{gl}_5(\bbQ\fps u)\fps z$. For a field $\Bbbk\supseteq\bbQ$ the proofs of \cref{lem:gau-comm,lem:gau-lift,lem:gau-scal,lem:gau-lead} hold verbatim over $\Bbbk$, so the solutions over $\Bbbk$ are $P^\natural\Bbbk\fps z$ by (b). If $P=P^\natural\cn$ has rational coefficients, the recursion of the surjectivity proof computes $\cn_n=(\coef{z^n}R^{(n)})_{00}(0)$ from rational data, so $\cn\in\bbQ\fps z$.
\end{proof}

\begin{lem}[The $\bbZ/5$-degree of the gauge equation]
\label{lem:gau-degree}
For $P=\sum_kz^kP_k$ with $P_k\in\mathfrak{gl}_5(\Bbbk\lser u)$ and $\delta\in\bbZ/5$ put $P^{[\delta]}:=\sum_kz^k\sum_{(a,b):\,k+b-a\equiv\delta}(P_k)_{ab}E_{ab}$, so that $P=\sum_\delta P^{[\delta]}$, and $P$ is graded if and only if $P=P^{[0]}$. Let $\operatorname{ge}(P):=zP'-AP+P\tilde A$. Then $\operatorname{ge}(P^{[\delta]})=\operatorname{ge}(P)^{[\delta+1]}$; in particular every degree part of a solution of the gauge equation is a solution.
\end{lem}

\begin{proof}
Give $z^kf(u)E_{ab}$ the degree $k+b-a$ modulo $5$. The derivative in $u$ does not change it and the factor $z$ adds $1$, so $zP'$ raises the degree of each homogeneous term by $1$. $A\,(fz^kE_{cd})=b_{c-1}fz^kE_{c-1,d}$ has degree $k+d-c+1$, and $(fz^kE_{cd})\tilde A=\tilde b_dfz^kE_{c,d+1}$ has degree $k+d+1-c$. So $\operatorname{ge}$ maps homogeneous series of degree $\delta$ to homogeneous series of degree $\delta+1$, and by additivity $\operatorname{ge}(P)^{[\delta+1]}=\operatorname{ge}(P^{[\delta]})$.
\end{proof}

\begin{prop}[Graded gauges]
\label{prop:gau-graded}
Let $A$, $\tilde A$ have the one-pole shape and satisfy (EP).
\begin{enumerate}
\item Graded regular gauges exist; they can be chosen rational if all $b_a$, $\tilde b_a$ are rational.
\item If $P$ is a graded regular gauge, the graded regular gauges are exactly the $P\cn(z)$ with $\cn\in\Bbbk\fps{z^5}^\times$.
\item For a graded regular gauge $P$, $P_{11}(0,z)=\sum_{m\geq0}(P_{5m})_{11}(0)\,z^{5m}$ is a unit of $\Bbbk\fps{z^5}$.
\item If $\zeta\in\Bbbk$, then $\operatorname{Rot}_iP=\sum_\delta\zeta^{-i\delta}P^{[\delta]}$; so $P$ is graded if and only if $\operatorname{Rot}_1P=P$, if and only if $\operatorname{Rot}_iP=P$ for all $i$, and $\frac15\sum_i\operatorname{Rot}_iP=P^{[0]}$. Moreover $A\,\zeta^{i(\mathrm{gr}+1)}=\zeta^i\zeta^{i(\mathrm{gr}+1)}A$, and the same for $\tilde A$; $\operatorname{Rot}_i$ maps solutions of $z\de_uY=AY$ to solutions, and $\operatorname{Rot}_i(PY)=(\operatorname{Rot}_iP)(\operatorname{Rot}_iY)$, $\operatorname{Rot}_i(Y\mathsf d(z))=\operatorname{Rot}_i(Y)\,\mathsf d(\zeta^{-i}z)$ for scalar series $\mathsf d$.
\end{enumerate}
\end{prop}

\begin{proof}
(i) Let $P^\natural$ be as in \cref{thm:gau-U}~(a), rational in the rational case by (d). Its $z^0$-term $\mathfrak d\Theta_U$ is diagonal, of degree $0$. By \cref{lem:gau-degree}, $P^{\mathrm{gr}}:=(P^\natural)^{[0]}$ solves the gauge equation, and $P^{\mathrm{gr}}_0=P^\natural_0$ is invertible. (ii) By \cref{thm:gau-U}~(b) the regular gauges are the $P\cn(z)$, $\cn\in\Bbbk\fps z^\times$. Write $\cn=\sum_{r\in\bbZ/5}\cn^{(r)}$ with $\cn^{(r)}:=\sum_{n\equiv r}\cn_nz^n$. As $P$ has degree $0$ and $z^n$ adds $n$ to the degree, $(P\cn)^{[r]}=P\cn^{(r)}$. So $P\cn$ is graded if and only if $P\cn^{(r)}=0$ for $r\neq0$, if and only if $\cn^{(r)}=0$ for $r\neq0$, because $\cn\mapsto P\cn$ is injective ($P=P^\natural\cn_P$ with $\cn_P\neq0$, and $\Bbbk\fps z$ is a domain); that is, if and only if $\cn\in\Bbbk\fps{z^5}$. (iii) For graded $P$ the entry $(P_k)_{11}$ vanishes unless $1\equiv1+k$, so $P_{11}(0,z)\in\Bbbk\fps{z^5}$, and its constant term is $(P_0)_{11}(0)=\cn(0)\mathfrak d(0)d^U_1(0)\neq0$ by \cref{thm:gau-U}~(b). (iv) The entry $(a,b)$ of $\coef{z^k}\operatorname{Rot}_iP$ is $\zeta^{i(a+1)-ik-i(b+1)}(P_k)_{ab}=\zeta^{-i(k+b-a)}(P_k)_{ab}$, which gives the first assertion; the equivalences follow because the decomposition $P=\sum_\delta P^{[\delta]}$ is direct, $\zeta^{-i\delta}\neq1$ for $i,\delta\not\equiv0$, and $\sum_i\zeta^{-i\delta}=5\,\delta_{\delta,0}$. The entry $(a,a+1)$ of $\zeta^{i(\mathrm{gr}+1)}A\zeta^{-i(\mathrm{gr}+1)}$ is $\zeta^{i(a+1)}b_a\zeta^{-i(a+2)}=\zeta^{-i}b_a$, which is the commutation rule. Hence, for a solution $Y$, $z\de_u(\operatorname{Rot}_iY)=\zeta^i\zeta^{i(\mathrm{gr}+1)}\bigl[(\zeta^{-i}z)\de_uY\bigr](u,\zeta^{-i}z)=\zeta^i\zeta^{i(\mathrm{gr}+1)}A\,Y(u,\zeta^{-i}z)=A\operatorname{Rot}_iY$. The two product rules are immediate from the definitions.
\end{proof}

\begin{rmk}[What the proof of Theorem~U uses; the role of the pole]
\label{rmk:gau-Uremarks}
The proof of \cref{thm:gau-U} uses only the one-pole shape, (EP) and the characteristic $0$ (in $\lambda_m(\tilde A^m)=5\neq0$, in the pivots $n+j/5\neq0$, and for $\mathfrak d$ and the constants of $\dd/\dd u$). The pole of $\tilde b_1$ enters exactly through the integrality statements, \cref{lem:gau-comm}~(f) and the window argument of \cref{lem:gau-lift}, and through the term $j/(5u)$ of $\chi^U_j$. For two systems without a pole (all ten weights units, with (EP)) the integrality statements hold trivially and $\chi^U_j\in\Bbbk\fps u$ for every $j$, so in the proof of \cref{lem:gau-scal}~(b) every equation $w_j'+\chi^U_jw_j=-\phi^Z_j$, $j=0,\dots,4$, has the solutions $\mathcal E_j\bigl(\mathsf a_j-\int_0^u\phi^Z_j/\mathcal E_j\,\dd u\bigr)$, $\mathsf a_j\in\Bbbk$, with $\mathcal E_j:=\exp(-\int_0^u\chi^U_j\,\dd u)$ (local letters). Then the proof of \cref{lem:gau-lead} gives the leading terms $\Theta_U\sum_j\mathsf a_j\mathcal E_j\mathcal Y^j$, the construction of \cref{thm:gau-U}~(a) started from $w^{(0)}:=\mathcal E_j\mathcal Y^j$ gives a solution $P^{(j)}$ for each $j$, and the proof of \cref{thm:gau-U}~(b), with the five leading terms in place of one, shows that $P^{(0)},\dots,P^{(4)}$ form a basis of $\Gau(A,\tilde A)$ over $\Bbbk\fps z$. So the solution module has rank $5$, and $P_{11}(0,z)=1$ no longer fixes a gauge (see \cref{sec:evi-gauge} for a computed example). This is not used.
\end{rmk}

% Section 3, part 5 of 7: (R1), the diagonal, the vanishing solutions, the gauge identity (G), the canonical gauge, symplecticity.

%-------------------------------------------------------------------------------
\subsection{The residue relation, the diagonal, the vanishing solutions and (G)}
\label{sec:gau-van}

From now on, unless said otherwise, $A$ is the quintic system of \cref{prop:gau-quintic} in a conifold frame and $\tilde A$ is the model pulled back along $v=v(u)$ (\cref{prop:gau-equalprod}). Both are self-dual cyclic systems with one simple pole over $\bbQ\fps u$, with (EP), $r=\tilde r=-1$ and $b_a(0)=\tilde b_a(0)$ for $a\neq1$, so \cref{thm:gau-U} applies over $\bbQ$.

\begin{lem}[The residue relation (R1)]
\label{lem:gau-R1}
Let $A,\tilde A$ be cyclic systems with one simple pole over $\Bbbk\fps u$, with residues $r,\tilde r$, and let $P\in\Gau(A,\tilde A)$; put $P(0,z):=\sum_kz^kP_k(0)\in\mathfrak{gl}_5(\Bbbk\fps z)$. Then $rE_{12}P(0,z)=\tilde r\,P(0,z)E_{12}$: column $1$ of $P(0,z)$ is $P_{11}(0,z)\mathbf e_1$, row $2$ is $P_{22}(0,z)\mathbf e_2^T$, and $rP_{22}(0,z)=\tilde rP_{11}(0,z)$. If $r=\tilde r$, then $P_{22}(0,z)=P_{11}(0,z)$.
\end{lem}

\begin{proof}
$A=u^{-1}rE_{12}+(\text{power series in }u)$, and likewise for $\tilde A$. The term $zP'$ has no $u^{-1}$, so the coefficient of $u^{-1}$ in $AP-P\tilde A$ vanishes: $rE_{12}P(0,z)=\tilde rP(0,z)E_{12}$. The entry $(a,b)$ reads $r\,\delta_{a1}P_{2b}(0,z)=\tilde r\,P_{a1}(0,z)\delta_{b2}$; for $a=1$, $b\neq2$ it gives $P_{2b}(0,z)=0$, for $a\neq1$, $b=2$ it gives $P_{a1}(0,z)=0$, and for $(a,b)=(1,2)$ it gives $rP_{22}(0,z)=\tilde rP_{11}(0,z)$.
\end{proof}

\begin{lem}[The diagonal of the gauge]
\label{lem:gau-diag}
Let $A,\tilde A$ be cyclic systems with one simple pole over $\Bbbk\fps u$ with (EP), and $P$ a regular gauge. Then $P_0=\diag(p_{0,0},\dots,p_{0,4})$, and $\prod_ap_{0,a}=\cn(0)^5\det\Theta_U(0)=p_{0,1}(0)^5\prod_a\rho^U_a$ is a constant, where $\cn(0)\in\Bbbk^\times$ is as in \cref{thm:gau-U}~(b) and $\rho^U_a:=d^U_a(0)/d^U_1(0)$, i.e.\ $\rho^U_1=1$, $\rho^U_a=\prod_{c=1}^{a-1}(\tilde b_c/b_c)(0)$ for $a=2,3,4$ and $\rho^U_0=\prod_{c=1}^4(\tilde b_c/b_c)(0)$, with $(\tilde b_1/b_1)(0):=\tilde r/r$. Assume the normalisation
\[
\text{(N0)}\qquad b_a(0)=\tilde b_a(0)\ \ (a\neq1),\qquad r=\tilde r .
\]
Then $p_{0,a}(0)=p_{0,1}(0)$ for all $a$; if moreover $p_{0,1}(0)=1$, as for the canonical gauge of \cref{def:gau-can}, then $p_{0,a}(0)=1$ for all $a$, $\prod_ap_{0,a}=1$, and $p_{0,0}^5=\prod_{c=0}^3(b_c/\tilde b_c)^{4-c}$. For the quintic system in a conifold frame and the pulled-back model, where (N0) holds, and a regular gauge with $p_{0,1}(0)=1$ (for instance $\Pc$), this gives
\[
p_{0,0}(u)=\tau(u)^{-1/4}/\varpi_0,\qquad\text{i.e.}\qquad \tau^{1/4}p_{0,0}=1/\varpi_0 .
\]
\end{lem}

\begin{proof}
By \cref{thm:gau-U}~(b), $P_0=\cn(0)\,\mathfrak d\,\Theta_U$ with $\cn(0)\neq0$; so $p_{0,a}=\cn(0)\mathfrak d\,d^U_a$ and $\prod_ap_{0,a}=\cn(0)^5\mathfrak d^5\det\Theta_U=\cn(0)^5\det\Theta_U(0)$, by the definition of $\mathfrak d$. At $u=0$, $p_{0,a}(0)=\cn(0)d^U_a(0)=p_{0,1}(0)\rho^U_a$. The recursion $d^U_{a+1}=d^U_a\tilde b_a/b_a$ at $u=0$, with $(\tilde b_1/b_1)(0)=\tilde r/r$ and, going round the cycle by (EP), $d^U_0=d^U_4\tilde b_4/b_4$, gives the formula for $\rho^U_a$. Under (N0) every ratio at $u=0$ is $1$, so $d^U_a(0)=1$ and $\rho^U_a=1$ for all $a$, and $p_{0,a}(0)=\cn(0)$. If $p_{0,1}(0)=1$, then $\cn(0)=1$, all $p_{0,a}(0)=1$, $\prod_ap_{0,a}=\det\Theta_U(0)=1$, and $p_{0,0}=\mathfrak d=(\det\Theta_U)^{-1/5}$. Since $d^U_a=\prod_{c<a}\tilde b_c/b_c$, each $c\in\{0,\dots,3\}$ occurs in $4-c$ of the $d^U_a$, so $\det\Theta_U=\prod_{c=0}^3(\tilde b_c/b_c)^{4-c}$.

For the two systems of \cref{prop:gau-quintic,prop:gau-equalprod} the ratios $b_c/\tilde b_c$, $c=0,\dots,3$, are
\[
\frac{\tc'}{v'},\qquad \frac{v}{u\,v'\varpi_0^2(1-u)^3\tc'^2},\qquad\frac{\tc'}{v'},\qquad\frac{\varpi_0}{(1-u)v'},
\]
so $p_{0,0}^5=\tc'^{\,4-6+2}\,v'^{-10}\,v^3u^{-3}\,\varpi_0^{-6+1}(1-u)^{-9-1}=v'^{-10}v^3u^{-3}\varpi_0^{-5}(1-u)^{-10}$. By (EP), $v'^5=v/(u(1-u)^5)$ (\cref{prop:gau-equalprod}~(ii)), hence $v'^{-10}=u^2(1-u)^{10}v^{-2}$ and $p_{0,0}^5=v/(u\varpi_0^5)=\tau^{-5/4}\varpi_0^{-5}$. Both $p_{0,0}$ and $\tau^{-1/4}/\varpi_0$ lie in $1+u\bbQ\fps u$, where fifth roots are unique, so they are equal.
\end{proof}

\begin{defn}[Vanishing solutions]
\label{def:gau-van}
\begin{enumerate}
\item For $m\in\bbZ$ let $\IC_m(v,z):=\Res_{p=0}p^m\re^{-\Wv(p)/z}\dd p$ be the formal residue: in $\re^{-\Wv(p)/z}=\re^{-p^4/(4z)}\re^{-v/(pz)}$ the coefficient of each power of $v$ is a Laurent series in $p$ with coefficients in $\bbQ[z^{-1}]$, and $\Res_{p=0}$ takes its coefficient of $p^{-1}$. Explicitly
\[
\IC_m(v,z)=\sum_{i\geq0,\ m+4i+1\geq0}\frac{(-1/(4z))^i}{i!}\,\frac{(-v/z)^{m+4i+1}}{(m+4i+1)!}\in\bbQ[z^{-1}]\fps v .
\]
The \emph{vanishing solution of the model} is $\YC:=(\IC_0,-\IC_{-1},-\IC_3,\IC_2,\IC_1)$, i.e.\ $\YC_b=\varepsilon_b\IC_{\nu_b}$ (\cref{def:set-model}).
\item For a cyclic system $A=\sum_ab_aE_{a,a+1}$ with one simple pole over $\bbQ\fps u$ put $y^{(0)}:=-\mathbf e_1$, $y^{(n+1)}_a:=\int_0^ub_a\,y^{(n)}_{a+1}\,\dd u$ (the integrands are power series by \cref{lem:gau-vanq}) and $Y^{\mathrm C}_A:=\sum_{n\geq0}z^{-n}y^{(n)}$. For the quintic system, $\YCQ:=Y^{\mathrm C}_A$ is the \emph{vanishing solution of the quintic}.
\end{enumerate}
\end{defn}

\begin{lem}[The model's vanishing solution]
\label{lem:gau-vanmod}
(i) $z\,\de_v\IC_m=-\IC_{m-1}$; (ii) $mz\IC_{m-1}-\IC_{m+3}+v\IC_{m-2}=0$. Hence $z\,\dd\YC_b/\dd v=b^{\mathrm m}_b\YC_{b+1}$, $b^{\mathrm m}=(1,-v^{-1},1,-1,-1)$, and $\YC(0,z)=-\mathbf e_1$.
\end{lem}

\begin{proof}
(i) Differentiate the residue coefficientwise in $v$, using $\de_v\re^{-\Wv/z}=-(pz)^{-1}\re^{-\Wv/z}$. (ii) We have $\frac{\dd}{\dd p}\bigl(p^m\re^{-\Wv/z}\bigr)=\bigl(mp^{m-1}-z^{-1}(p^{m+3}-vp^{m-2})\bigr)\re^{-\Wv/z}$, since $\Wv'(p)=p^3-vp^{-2}$, and the residue of the derivative of a Laurent series vanishes. With $\YC=(\IC_0,-\IC_{-1},-\IC_3,\IC_2,\IC_1)$, (i) gives $z\de_v\YC_0=-\IC_{-1}=\YC_1$, $z\de_v\YC_2=\IC_2=\YC_3$, $z\de_v\YC_3=-\IC_1=-\YC_4$, $z\de_v\YC_4=-\IC_0=-\YC_0$, and $z\de_v\YC_1=\IC_{-2}=v^{-1}\IC_3=-v^{-1}\YC_2$ by (ii) at $m=0$. At $v=0$ only the terms with $m+4i+1=0$ survive; among $m\in\{0,-1,3,2,1\}$ this happens only for $m=-1$, $i=0$, with value $1$; so $\YC(0,z)=-\mathbf e_1$.
\end{proof}

\begin{lem}[Shape of $\YC$]
\label{lem:gau-vanshape}
For $n\geq0$ and $b\in\bbZ/5$, with $m:=\lfloor n/5\rfloor$,
\[
\coef{z^{-n}}\YC_b=\varepsilon_b\,(-1)^n\,\frac{v^{n-m}}{4^m\,m!\,(n-m)!}\quad\text{if }n\equiv\nu_b+1\pmod5,\qquad \coef{z^{-n}}\YC_b=0\ \text{otherwise}.
\]
Hence $\YC_b(v,z)=\varsigma^{\nu_b+1}y_b(z/\varsigma^4)$ with $y_b(t):=\YC_b(1,t)$ and $\varsigma=v^{1/5}$; the lowest terms are $\YC=(-v/z,\,-1,\,-v^4/(24z^4),\,-v^3/(6z^3),\,v^2/(2z^2))+\dots$; and at $v=v(u)$ the coefficient of $z^{-n}$ is $O(u^{n-m})$, so $\YC(v(u),z)\in\bbQ[z^{-1}]\fps u^5$ and $\YC(v(u),z)$ has no positive power of $z$.
\end{lem}

\begin{proof}
In \cref{def:gau-van}~(i) with $m'=\nu_b$, the term of index $i$ is a multiple of $v^{m'+4i+1}z^{-n}$ with $n=m'+5i+1$; as $\nu_b+1\in\{0,\dots,4\}$, this gives $n\equiv\nu_b+1$ and $i=\lfloor n/5\rfloor=m$, the exponent of $v$ is $n-m$, the sign is $(-1)^{i+m'+4i+1}=(-1)^n$, and the denominator is $4^mm!\,(n-m)!$. Since $v^{n-m}z^{-n}=\varsigma^{n-5m}(z/\varsigma^4)^{-n}$ and $n-5m=\nu_b+1$, the homogeneity follows. For $v=u+O(u^2)$ the coefficient of $z^{-n}$ is $O(u^{n-m})$ with $n-m\geq4n/5$, so each power of $u$ receives contributions from finitely many $n$.
\end{proof}

\begin{lem}[Uniqueness of regular solutions]
\label{lem:gau-uniq}
Let $B=u^{-1}B_{-1}+\sum_{j\geq0}u^jB_j$ with $B_j\in\mathfrak{gl}_5(\Bbbk)$ independent of $z$. If $Y\in\Bbbk\lser z\fps u^5$ satisfies $z\de_uY=BY$ and $Y(0)=0$, then $Y=0$.
\end{lem}

\begin{proof}
Write $Y=\sum_nu^nY^{(n)}$. The coefficient of $u^n$ in $z\de_uY=BY$ gives $\bigl((n+1)z-B_{-1}\bigr)Y^{(n+1)}=\sum_{j=0}^nB_jY^{(n-j)}$. The determinant $\det((n+1)z\,1-B_{-1})$ is a polynomial of degree $5$ in $z$ with leading coefficient $(n+1)^5\neq0$, hence a unit of $\Bbbk\lser z$ (for $B_{-1}=rE_{12}$ it is $((n+1)z)^5$). By induction all $Y^{(n)}$ vanish.
\end{proof}

A cyclic system with one simple pole is of this form ($B_{-1}=rE_{12}$), so its solutions in $\bbQ\lser z\fps u^5$ are determined by their values at $u=0$. By \cref{lem:gau-vanmod} and the chain rule, $\YC(v(u),z)$ is the solution of $z\de_uY=\tilde AY$ with value $-\mathbf e_1$, and by the next lemma $\YCQ$ is that of the quintic system; the vanishing solutions are characterised by this value, not by the absence of positive powers of $z$.

\begin{lem}[The quintic's vanishing solution]
\label{lem:gau-vanq}
Let $A=\sum_ab_aE_{a,a+1}$ be a cyclic system with one simple pole over $\bbQ\fps u$, with residue $r$.
\begin{enumerate}
\item[(a)] $y^{(n)}_a=0$ unless $n\equiv1-a\pmod5$, and then $\ord_uy^{(n)}_a=(4n+\nu_a+1)/5$ exactly. If $b_a(0)=\tilde b_a(0)$ for $a\neq1$ and $r=\tilde r$, where $\tilde A$ is the pulled-back model, the leading coefficient of $y^{(n)}_a$ equals that of $\coef{z^{-n}}\YC_a(v(u),z)$.
\item[(b)] Every integrand $b_ay^{(n)}_{a+1}$ is a power series: the pole of $b_1$ meets $y^{(n)}_2$, which is non-zero only for $n\equiv4\pmod5$ and then has order $(4n+4)/5\geq4$. In particular no logarithm arises.
\item[(c)] $Y^{\mathrm C}_A\in\bbQ[z^{-1}]\fps u^5$, the coefficient of $z^{-n}$ having order at least $4n/5$; it solves $z\de_uY=AY$, and $Y^{\mathrm C}_A(0)=-\mathbf e_1$.
\item[(d)] $\coef{z^{-1}}Y^{\mathrm C}_A=-\bigl(\int_0^ub_0\,\dd u\bigr)\mathbf e_0$ and $\coef{u^0}(Y^{\mathrm C}_A)_0=0$. For the quintic system in a conifold frame, $\coef{z^{-1}}\YCQ=-\tc\,\mathbf e_0$.
\end{enumerate}
\end{lem}

\begin{proof}
Put $e(n,a):=(4n+\nu_a+1)/5$; for $n\equiv1-a$ it is a non-negative integer, as $\nu_a\equiv-a$ and $\nu_a+1\geq0$. Write $b_a=b^\circ_au^{-\delta_{a1}}(1+O(u))$ with $b^\circ_a:=b_a(0)$ for $a\neq1$ and $b^\circ_1:=r$ (local letters). We show by induction on $n$ that for $n\equiv1-a$, $y^{(n)}_a=c^{(n)}_au^{e(n,a)}(1+O(u))$ with $c^{(n)}_a\neq0$, and $y^{(n)}_a=0$ otherwise. For $n=0$, $y^{(0)}=-\mathbf e_1$ and $e(0,1)=0$. The component $y^{(n+1)}_a$ is built from $y^{(n)}_{a+1}$, which vanishes unless $n\equiv1-(a+1)$, i.e.\ $n+1\equiv1-a$. In that case, using $\nu_{a+1}=\nu_a-1+5\delta_{a1}$, the integrand is $b_ay^{(n)}_{a+1}=b^\circ_ac^{(n)}_{a+1}u^{e(n+1,a)-1}(1+O(u))$ with $e(n+1,a)-1=(4n+\nu_a)/5\geq0$, a power series, and integration gives $c^{(n+1)}_a=b^\circ_ac^{(n)}_{a+1}/e(n+1,a)\neq0$. This proves the first part of (a) and (b); for $a=1$ the integrand involves $y^{(n)}_2$ with $n\equiv4$, of order $e(n,2)=(4n+4)/5\geq4$. The pulled-back model satisfies the same recursion for the leading coefficients with $\tilde b^\circ_a$ in place of $b^\circ_a$, so under the stated hypothesis the leading coefficients of $Y^{\mathrm C}_A$ and $Y^{\mathrm C}_{\tilde A}$ agree; and $Y^{\mathrm C}_{\tilde A}=\YC(v(u),z)$ by \cref{lem:gau-uniq}, both being solutions of $z\de_uY=\tilde AY$ with value $-\mathbf e_1$ (the first by (c) below, applied to $\tilde A$). (c) The order bound is $e(n,a)\geq4n/5$, so each power of $u$ receives finitely many $n$. The definition says $(y^{(n+1)})'=Ay^{(n)}$, hence $z\de_uY^{\mathrm C}_A=\sum_{n\geq1}z^{1-n}Ay^{(n-1)}=AY^{\mathrm C}_A$; and $y^{(n)}(0)=0$ for $n\geq1$. (d) $y^{(1)}_a=-\delta_{a0}\int_0^ub_0$, and $(Y^{\mathrm C}_A)_0(0)=y^{(0)}_0=0$. For the quintic, $b_0=\tc'$ and $\tc(0)=0$.
\end{proof}

In a conifold frame the hypothesis of (a) holds (\cref{prop:gau-equalprod}).

\begin{thm}[The gauge identity (G)]
\label{thm:gau-G}
For every $P\in\Gau(A,\tilde A)$, in particular for every regular gauge from the pulled-back model to the quintic system,
\begin{equation}
\label{eq:gau-G}
P(u,z)\,\YC(v(u),z)=P_{11}(0,z)\,\YCQ(u,z)\qquad\text{in }\bbQ\lser z\fps u^5 .
\end{equation}
\begin{enumerate}
\item[(a)] For a graded regular gauge $P$, $P_{11}(0,z)$ is a unit of $\bbQ\fps{z^5}$, $\Pc:=P\,P_{11}(0,z)^{-1}$ is a graded regular gauge, and $\Pc$ is the unique regular gauge $P'$ with $P'\,\YC(v(u),z)=\YCQ$.
\item[(b)] For a regular gauge $P$, $P\,\YC(v(u),z)$ has no positive power of $z$ if and only if $P_{11}(0,z)$ is a constant; the regular gauges with this property are the $\mu\Pc$, $\mu\in\bbQ^\times$.
\item[(c)] For any two cyclic systems $A,\tilde A$ with one simple pole over $\bbQ\fps u$ and every $P\in\Gau(A,\tilde A)$: $P\,Y^{\mathrm C}_{\tilde A}=P_{11}(0,z)\,Y^{\mathrm C}_A$.
\end{enumerate}
\end{thm}

\begin{proof}
We prove (c), of which \eqref{eq:gau-G} is the case of the two systems of this subsection: $\YC(v(u),z)=Y^{\mathrm C}_{\tilde A}$ by \cref{lem:gau-uniq}, both being solutions of $z\de_uY=\tilde AY$ with value $-\mathbf e_1$ (\cref{lem:gau-vanmod,lem:gau-vanq}~(c)). Put $\tilde Y:=Y^{\mathrm C}_{\tilde A}\in\bbQ[z^{-1}]\fps u^5$. The product $P\tilde Y$ lies in $\bbQ\lser z\fps u^5$: at each order in $u$ it is a finite sum of products of a power series in $z$ and a polynomial in $z^{-1}$. It solves the first system: $z\de_u(P\tilde Y)=(AP-P\tilde A)\tilde Y+P\tilde A\tilde Y=A(P\tilde Y)$. Its value at $u=0$ is $P(0,z)(-\mathbf e_1)=-P_{11}(0,z)\mathbf e_1$ by \cref{lem:gau-R1}. $P_{11}(0,z)Y^{\mathrm C}_A$ is also a solution, $P_{11}(0,z)$ being constant in $u$, with the same value by \cref{lem:gau-vanq}~(c). \Cref{lem:gau-uniq} with $B=A$ gives the identity.

(a) $P_{11}(0,z)$ is a unit of $\bbQ\fps{z^5}$ by \cref{prop:gau-graded}~(iii), so $\Pc$ is a graded regular gauge by \cref{prop:gau-graded}~(ii), and $(\Pc)_{11}(0,z)=1$; by \eqref{eq:gau-G}, $\Pc\YC(v(u))=\YCQ$. Conversely, let $P'$ be any regular gauge with $P'\YC(v(u))=\YCQ$. By \cref{thm:gau-U}~(b), (d), $P'=\Pc\cn(z)$ for some $\cn\in\bbQ\fps z^\times$, so that $\YCQ=\cn(z)\Pc\YC(v(u))=\cn(z)\YCQ$, i.e.\ $(\cn(z)-1)\YCQ=0$ in $\bbQ\lser z\fps u^5$; the coefficient of $u^0$ is $-(\cn(z)-1)\mathbf e_1$, so $\cn=1$. (b) If $P_{11}(0,z)=\mu$ is constant, then $P\YC(v(u))=\mu\YCQ$ lies in $\bbQ[z^{-1}]\fps u^5$. Conversely, the coefficient of $u^0$ of $P\YC(v(u))$ is $-P_{11}(0,z)\mathbf e_1$ with $P_{11}(0,z)\in\bbQ\fps z$, so it has no positive power of $z$ only if $P_{11}(0,z)$ is constant. With $P=\Pc\cn(z)$ we have $P_{11}(0,z)=\cn(z)$, which gives the last assertion.
\end{proof}

\begin{defn}[The canonical gauge]
\label{def:gau-can}
In a conifold frame, the \emph{canonical gauge} $\Pc$ is the regular gauge from the pulled-back model to the quintic system with $(\Pc)_{11}(0,z)=1$. It exists and is unique: $\Pc=P\,P_{11}(0,z)^{-1}$ for any graded regular gauge $P$ (\cref{thm:gau-G}~(a)), two graded regular gauges differ by a unit of $\bbQ\fps{z^5}$ (\cref{prop:gau-graded}), and a regular gauge $P'=\Pc\cn(z)$ (\cref{thm:gau-U}~(b)) with $P'_{11}(0,z)=1$ has $\cn=1$. Equivalently, $\Pc$ is the unique regular gauge with $\Pc\YC(v(u))=\YCQ$. It is graded, rational (\cref{thm:gau-U}~(d)), colour covariant, $\operatorname{Rot}_i\Pc=\Pc$ (\cref{prop:gau-graded}~(iv)), and $p_{0,a}(0)=1$ for its diagonal $z^0$-term (\cref{lem:gau-diag}). We write $\Pc=\sum_kz^k\sum_bp_{k,b}(u)E_{b+k,b}$.
\end{defn}

%-------------------------------------------------------------------------------
\subsection{Symplecticity}
\label{sec:gau-symp}

\begin{prop}[Symplecticity]
\label{prop:gau-symp}
Let $A,\tilde A$ be self-dual cyclic systems with one simple pole over $\Bbbk\fps u$ with (EP) and $r=\tilde r$, and let $P$ be a regular gauge from $\tilde A$ to $A$, $\cn(z):=P_{11}(0,z)$. Then
\begin{equation}
\label{eq:gau-symp}
P(u,-z)^T\eta\,P(u,z)=\cn(z)\cn(-z)\,\eta .
\end{equation}
In particular, for the two systems of \cref{sec:gau-van}, $\Pc$ is symplectic, $\Pc(u,-z)^T\eta\Pc(u,z)=\eta$; the symplectic regular gauges are the $\Pc\cn(z)$ with $\cn(z)\cn(-z)=1$, so symplecticity constrains only the even part of $\log(\cn/\cn(0))$; for graded gauges, $\cn=\pm\exp(a_1z^5+a_3z^{15}+a_5z^{25}+\cdots)$. Among the symplectic regular gauges, only (G) singles out $\Pc$.
\end{prop}

\begin{proof}
Put $S:=P(u,-z)^T\eta P(u,z)$. Replacing $z$ by $-z$ in the gauge equation gives $zP'(-z)=-AP(-z)+P(-z)\tilde A$, so $zP'(-z)^T=-P(-z)^TA^T+\tilde A^TP(-z)^T$. Hence
\[
zS'=\bigl(-P(-z)^TA^T+\tilde A^TP(-z)^T\bigr)\eta P(z)+P(-z)^T\eta\bigl(AP(z)-P(z)\tilde A\bigr)=\tilde A^TS-S\tilde A,
\]
because $A^T\eta=\eta A$ by self-duality (\cref{def:gau-onepole}). With $T:=\eta^{-1}S$ and $\tilde A^T\eta=\eta\tilde A$ this is $zT'=\tilde AT-T\tilde A$, and $T$ has coefficients in $\mathfrak{gl}_5(\Bbbk\fps u)$. By \cref{thm:gau-U}~(b) for the pair $(\tilde A,\tilde A)$, $T=t(z)\,1$ for some $t\in\Bbbk\fps z$; no invertibility of $T_0$ is needed. Compare the entries $(2,1)$ at $u=0$: by \cref{lem:gau-R1} column $1$ of $P(0,z)$ is $\cn(z)\mathbf e_1$, so $S_{21}(0,z)=\sum_aP_{a2}(0,-z)\eta_{a1}\cn(z)=5P_{22}(0,-z)\cn(z)=5\cn(-z)\cn(z)$, using $\eta_{a1}=5\delta_{a2}$ and $P_{22}(0,\cdot)=P_{11}(0,\cdot)$ ($r=\tilde r$); on the other hand $S_{21}=t(z)\eta_{21}=5t(z)$. So $t(z)=\cn(z)\cn(-z)$. For the consequences: $(\Pc)_{11}(0,z)=1$; every regular gauge is $\Pc\cn(z)$ with $(\Pc\cn)_{11}(0,z)=\cn(z)$; $\cn(z)\cn(-z)=1$ forces $\cn(0)=\pm1$ and is equivalent to the oddness of $\log(\cn/\cn(0))$; for graded gauges $\cn\in\Bbbk\fps{z^5}$.
\end{proof}

\begin{cor}[The dual identity]
\label{cor:gau-Gstar}
As row vectors in $\bbQ\lser z\fps u$,
\[
\YCQ(u,-z)^T\eta\,\Pc(u,z)=\YC(v(u),-z)^T\eta .
\]
More generally, $\YCQ(-z)^T\eta P(z)=P_{11}(0,z)\,\YC(-z)^T\eta$ for every regular gauge $P$.
\end{cor}

\begin{proof}
For a regular gauge, $P_{11}(0,z)$ is a unit of $\bbQ\fps z$, its constant term being $(P_0)_{11}(0)\neq0$. Substituting $-z$ for $z$ (a ring automorphism of $\bbQ\lser z\fps u$) in \eqref{eq:gau-G} gives $\YCQ(-z)=P_{11}(0,-z)^{-1}P(-z)\YC(-z)$. Transposing, multiplying on the right by $\eta P(z)$ and using \eqref{eq:gau-symp},
\[
\YCQ(-z)^T\eta P(z)=P_{11}(0,-z)^{-1}P_{11}(0,z)P_{11}(0,-z)\,\YC(-z)^T\eta=P_{11}(0,z)\,\YC(-z)^T\eta .\qedhere
\]
\end{proof}

\input{sections/s03_gauge_6}



